\chapter{Quantization and supported transfer}
\label{cyquant:chapter}

A classical contraction identifies a field complex with a smaller complex.
Quantization asks for more: the interaction, the measure, and the collision
operations must descend with their identities.  Three explicit constructions
will distinguish these requirements.  On $K3\times E$, harmonic forms give a
strict model of the classical holomorphic Chern--Simons Lie algebra, while
its full residual action retains two cubic terms.  In a nodal field algebra,
the BRST differential retains both the coordinate relation and its momentum
fields.  In a rational current algebra, a completed product and its curved
differential can be constructed through specified powers of the deformation
parameter.  Each construction has its own carrier and its own comparison
problem.

The final example concerns support.  Dividing by a normal coordinate on
principal parts has a right inverse with a nonzero contact term.  That term
survives in cohomology and prevents a particular division homotopy from
becoming module-linear.  Its explicit representative makes it possible to
state the remaining relative comparison equation, including the normal
derivatives and the vacuum.

\section{The full harmonic Chern--Simons model}
\label{cyquant:sec:residual}

Let $X=S\times E$, where $S$ is a compact K3 surface and $E$ is an
elliptic curve.  Fix a product Ricci-flat K\"ahler metric, a nowhere vanishing
holomorphic three-form $\Omega$, and a finite-dimensional complex Lie
algebra $\mathfrak g$ with a nondegenerate invariant symmetric form
$\kappa$.  We work about the trivial gauge background.  The differential
graded Lie algebra and its cyclic pairing are
\begin{equation}\label{cyquant:eq:hcs}
 L=\mathfrak g\otimes\Omega^{0,\bullet}(X),\qquad
 d=\bar\partial,\qquad
 \pi(x\alpha,y\beta)=\int_X\Omega\wedge\alpha\wedge\beta\,
                         \kappa(x,y).
\end{equation}
The pairing has degree $-3$.  The shifted space of BV fields is $L[1]$;
in particular it includes all four Dolbeault degrees.  Its odd symplectic
pairing is
\[
 \sigma(sx,sy)=(-1)^{|x|}\pi(x,y),\qquad |s|=-1.
\]
The restriction to Dolbeault degree one is the classical physical field
space.  The full BV extension also contains ghosts and antifields
\cite[\S1.5]{cyquant:CostelloLi}.

Choose parallel harmonic representatives $\epsilon$ and $u$ of the
elliptic $(0,1)$-class and the K3 $(0,2)$-class.  Normalize them so that
$\int_X\Omega\wedge\epsilon u=1$.  Their algebra is
\begin{equation}\label{cyquant:eq:harmonic-algebra}
 H=\mathbb C[u]/(u^2)\otimes\Lambda(\epsilon),
 \qquad |u|=2,\quad |\epsilon|=1.
\end{equation}
Write $\tau(\epsilon u)=1$, with $\tau$ zero in the other degrees.

\begin{proposition}\label{cyquant:prop:strict-harmonic}
The inclusion $\mathfrak g\otimes H\longrightarrow L$ is a cyclic
dg Lie quasi-isomorphism.  Homotopy transfer along the product Hodge
contraction has no operations of arity at least three on
$\mathfrak g\otimes H$.  Its binary bracket is
\[
 [x\otimes v,y\otimes w]=[x,y]\otimes vw.
\]
\end{proposition}

\begin{proof}
The Hodge dimensions of $S$ and $E$, together with K\"unneth, give the
four classes $1,\epsilon,u,\epsilon u$.  Their chosen representatives are
parallel and closed under multiplication, so their inclusion is a
multiplicative quasi-isomorphism.  The trace normalization identifies its
pairing with $\kappa\otimes\tau$.

The Hodge homotopy kills harmonic forms.  Every rooted binary tree with
at least three leaves has an internal edge immediately above a vertex
whose two inputs are harmonic.  Their bracket is harmonic, and the
homotopy on that edge kills it.  Thus every higher transferred tree is
zero.  The binary tree has no such internal edge and gives the displayed
bracket.
\end{proof}

The strictness of this Lie model does not remove its cubic Hamiltonian.
Write the universal unshifted field, with its coordinate coefficients on
the right, as
\[
 \omega=1c+\epsilon a+ub+\epsilon u d.
\]
The coordinate degrees are $|c|=1$, $|a|=0$, $|b|=-1$, and $|d|=-2$.
Consequently $c,b$ are odd and $a,d$ are even.  The residual action is the
degree-zero function
\begin{equation}\label{cyquant:eq:residual-action}
 S_H=\frac16\tau\kappa(\omega,[\omega,\omega])
     =\frac12\kappa(d,[c,c])-\kappa(c,[a,b]).
\end{equation}
To compute it, a nonzero trace must have form degree three.  The possible
patterns are $(0,0,3)$, $(0,1,2)$, and $(1,1,1)$.  The last vanishes
because $\epsilon^2=0$.  Moving the right-hand coordinate coefficients
past the forms in the first two patterns gives exactly the two signs in
\eqref{cyquant:eq:residual-action}.

For example, take $\mathfrak g=\mathfrak{sl}_2$ with
$e=E_{12}$, $f=E_{21}$, $h=\operatorname{diag}(1,-1)$ and
$\kappa(x,y)=\operatorname{tr}(xy)$.  Since
$[e,f]=h$, $\kappa(h,h)=2$, and $\kappa(e,f)=1$, the coefficients of
$c_h a_e b_f$ and $c_e c_f d_h$ in $S_H$ are respectively $-2$ and $2$.
The action vanishes on the restriction $c=0$, but it is nonzero on the
full residual field space.

\begin{proposition}\label{cyquant:prop:finite-qme}
Use left graded derivatives and the constant-density BV Laplacian
\[
 \Delta=\sum_{i,j}\kappa^{ij}
 \left(\frac{\partial}{\partial c_i}\frac{\partial}{\partial d_j}
       +\frac{\partial}{\partial a_i}\frac{\partial}{\partial b_j}\right).
\]
For the bracket $\{-,-\}$ generated by $\Delta$, one has
$\{S_H,S_H\}=0$ and $\Delta S_H=0$.  The Hamiltonian differential is
\begin{equation}\label{cyquant:eq:residual-Q}
 Qc=\tfrac12[c,c],\qquad Qa=[c,a],\qquad Qb=[c,b],\qquad
 Qd=[c,d]-[a,b].
\end{equation}
Thus $S_H$ solves the finite algebraic quantum master equation.
\end{proposition}

\begin{proof}
Expand $d_{\mathrm{tot}}\omega=\frac12[\omega,\omega]$, where
$d_{\mathrm{tot}}=1_H\otimes Q$ uses the tensor sign.
Its four coefficients give
\eqref{cyquant:eq:residual-Q}.  Invariance of $\kappa$ identifies this
differential with the Hamiltonian derivation of $S_H$.  The graded Jacobi
identity gives $Q^2=0$ and the classical master equation.  A metric Lie
algebra in characteristic zero is unimodular: each adjoint endomorphism
is skew with respect to its nondegenerate symmetric form, hence has trace
zero.  The contractions in $\Delta S_H$ are precisely these adjoint
traces, so they vanish.
\end{proof}

This finite cyclic construction is the algebraic residual mechanism of
\cite[\S3.1]{cyquant:CattaneoMnev}.  Its differential on $H$ is zero,
so the quadratic kinetic term restricts to zero.  The nondegenerate odd
symplectic form supplies the BV bracket, not a nondegenerate quadratic
action.  An effective quantum action for the analytic theory further
requires a regulated integration over nonzero modes, its local
counterterms, and a choice of residual BV integration.  Neither the
finite master equation nor $\chi(S\times E)=0$ computes those data.

\section{A product propagator with one time variable}
\label{cyquant:sec:propagator}

Let $A$ and $B$ be closed Hermitian manifolds.  On their Dolbeault
complexes write $d_A=\bar\partial_A$, $q_A=d_A^*$,
$\Delta_A=d_Aq_A+q_Ad_A$, and $\Gamma_A=(-1)^{\deg}$, and use analogous
notation for $B$.  All adjoints refer to the specified metrics.  On the
product complex,
\begin{equation}\label{cyquant:eq:product-operators}
 \begin{aligned}
 d&=d_A\otimes1+\Gamma_A\otimes d_B,\\
 q&=q_A\otimes1+\Gamma_A\otimes q_B,\\
 \Delta&=\Delta_A\otimes1+1\otimes\Delta_B.
 \end{aligned}
\end{equation}
The last identity follows because $\Gamma_A$ anticommutes with both odd
operators.  If $K_A(t)=e^{-t\Delta_A}$, then
$K(t)=K_A(t)\otimes K_B(t)$.

\begin{proposition}\label{cyquant:prop:product-propagator}
For $0<\varepsilon<L$, the product Hodge propagator is
\begin{equation}\label{cyquant:eq:propagator}
 P_{[\varepsilon,L]}=
 \int_{\varepsilon}^{L}
 \bigl(q_AK_A(t)\otimes K_B(t)
       +\Gamma_AK_A(t)\otimes q_BK_B(t)\bigr)\,dt.
\end{equation}
It satisfies
\[
 dP_{[\varepsilon,L]}+P_{[\varepsilon,L]}d=K(\varepsilon)-K(L).
\]
Let $H_A,H_B$ be the orthogonal harmonic projections, put
$H=H_A\otimes H_B$, and let $G$ be the inverse of $\Delta$ on $H^\perp$,
extended by zero on $H$.  Then $P=qG$ satisfies $dP+Pd=1-H$.
If $\lambda_1>0$ is the first positive eigenvalue of $\Delta$, then
\begin{equation}\label{cyquant:eq:propagator-bound}
 \|P-P_{[\varepsilon,L]}\|_{L^2\to L^2}
 \leq \sqrt{\varepsilon}+
             \frac{e^{-L\lambda_1}}{\sqrt{\lambda_1}}.
\end{equation}
\end{proposition}

\begin{proof}
The integrand is $qK(t)$.  Since $dq+qd=\Delta$, its anticommutator
with $d$ is $-\partial_tK(t)$, proving the finite-cutoff identity.
On a positive $\lambda$-eigenspace, $\|q\|\leq\sqrt\lambda$.
The omitted short-time integral is bounded there by
$(1-e^{-\varepsilon\lambda})/\sqrt\lambda\leq\sqrt\varepsilon$,
using $1-e^{-x}\leq\min(x,1)$.  The long-time tail is bounded by
$e^{-L\lambda}/\sqrt\lambda$.  Taking the supremum proves
\eqref{cyquant:eq:propagator-bound}.  On $H$, $qH=0$, so there is no
term growing linearly with $L$.  The spectral theorem gives the limit
$qG$, and elliptic regularity gives its action on smooth forms.
\end{proof}

There is only one time integral in \eqref{cyquant:eq:propagator}.
On a joint eigenspace with eigenvalues $\lambda,\mu>0$, its Green
denominator is $\lambda+\mu$.  Multiplying two separate Green operators
would instead give $\lambda\mu$.  The mixed harmonic sectors also
remain:
\[
 P|_{\operatorname{im}H_A\otimes\ker H_B}=\Gamma_AH_A\otimes P_B,
 \qquad P|_{\ker H_A\otimes\operatorname{im}H_B}=P_A\otimes H_B.
\]
The sector harmonic in both factors contributes zero.  If $h_A$ is an even
harmonic vector and $d_Bc=bd$ with $b\ne0$, then
$P(h_A\otimes d)=b^{-1}h_A\otimes c$.
The sequential contraction
$P_A\otimes1+\Gamma_AH_A\otimes P_B$ is another valid contraction;
it is a different choice from the product-metric Hodge contraction.

On $E=\mathbb C/(\mathbb Z+\tau\mathbb Z)$, write
$z=u+\tau v$, $\tau=a+ib$, $b>0$, and give the metric
$g=(V/b)|dz|^2$ area $V$.  In the convention
$\Delta_{\bar\partial}=\Delta_{\mathrm{real}}/2$, the characters and
eigenvalues are
\[
 e_{mn}(u,v)=e^{2\pi i(mu+nv)},\qquad
 \lambda_{mn}=\frac{2\pi^2|m\tau-n|^2}{Vb}.
\]
The scalar heat kernel is therefore
\begin{equation}\label{cyquant:eq:elliptic-heat}
 K_E(t;z,z')=\frac1V\sum_{m,n\in\mathbb Z}
       e^{-t\lambda_{mn}}e_{mn}(z)\overline{e_{mn}(z')}.
\end{equation}
Every summand is periodic in both variables.  Parallel orthonormal form
frames give the degree-one formula.  These conventions agree with the
Dolbeault Hodge identities in \cite[Ch.~VI, \S5]{cyquant:Demailly}.

For fixed positive cutoffs the propagator kernel is smooth even on the
diagonal.  Its limiting singularity has a different order from the scalar
Green kernel.  In real dimension six the positive flat Laplacian has
\[
 G(r)=\frac1{4\pi^3r^4},\qquad G'(r)=-\frac1{\pi^3r^5}.
\]
The first-order adjoint thus produces leading coefficients of order
$r^{-5}$ in fixed orthonormal frames.  Directly,
\[
 \int_0^\infty r t^{-4}e^{-r^2/(4t)}\,dt=\frac{128}{r^5},
 \qquad
 \int_\varepsilon^L r t^{-4}e^{-r^2/(4t)}\,dt
 \leq\frac r3(\varepsilon^{-3}-L^{-3}).
\]
The first expression is locally integrable in six dimensions; the second
has no short-distance pole.  These linear identities underlie the
heat-kernel construction of \cite[\S\S7.2.3--7.2.5]{CG2}.  An interacting
master equation requires the corresponding local functional calculation.

\section{Cyclic transfer and a nonformal minimal algebra}
\label{cyquant:sec:cyclic}

Let $(A,d,m_2)$ be a dg algebra with a cyclic pairing, and let
$(H,0)\mathrel{\substack{\xrightarrow{\ i\ }\\[-.7ex]
                         \xleftarrow[\ p\ ]{}}}(A,d)$
be a contraction with a homotopy $h$ of degree $-1$:
\[
 pi=1,\qquad dh+hd=1-ip,\qquad h^2=hi=ph=0.
\]
A cyclic transfer also requires $p,i,h$ to be compatible with the
pairing in the shifted sign convention.  The transferred operations are
the signed planar-tree compositions of $i,m_2,h,p$: leaves carry $i$,
vertices $m_2$, internal edges $h$, and the root $p$.  This gives
$|m_n|=2-n$.  It preserves cyclicity when the contraction is cyclic
\cite{cycmp:HTT}.  A linear contraction without these pairing conditions
does not by itself supply cyclic transfer.

We use the following unsuspended convention throughout this section:
\begin{equation}\label{cyquant:eq:stasheff-sign}
 \sum_{r+s+t=n}
 (-1)^{r+st+(2-s)(|x_1|+\cdots+|x_r|)}
 m_{r+1+t}(x_1,\ldots,x_r,m_s(x_{r+1},\ldots,x_{r+s}),
                         x_{r+s+1},\ldots,x_n)=0.
\end{equation}
For a degree-$-3$ pairing, cyclicity means
\begin{equation}\label{cyquant:eq:cyclic-sign}
 \begin{aligned}
 \langle m_n(x_1,\ldots,x_n),x_{n+1}\rangle
  ={}&(-1)^{n+|x_1|(|x_2|+\cdots+|x_{n+1}|)}\\
 &\hspace{-3em}\cdot
 \langle m_n(x_2,\ldots,x_{n+1}),x_1\rangle.
 \end{aligned}
\end{equation}
Both sides place the operation in the first pairing slot.  The shift,
the argument order, and the sign must be transported together when using
another convention \cite[\S2]{cyquant:Cho}.

\begin{proposition}\label{cyquant:prop:quartic-model}
Over a characteristic-zero field, let $A_\lambda$ have basis $1,a,b,v$
in degrees $0,1,2,3$.  Give it a strict unit $1$, the nonzero products
\[
 m_2(a,b)=m_2(b,a)=v,\qquad m_3(a,a,a)=\lambda b,
\]
and no other products on nonunit inputs.  Give it the pairing
\[
 \langle1,v\rangle=\langle v,1\rangle
 =\langle a,b\rangle=\langle b,a\rangle=1,
\]
zero on the other basis pairs.  This is a minimal cyclic
$A_\infty$-algebra of dimension three.  If $\lambda\ne0$, it is
nonformal and has both a nonzero ternary product and a nonzero binary
product on its augmentation ideal.
\end{proposition}

\begin{proof}
The binary product is associative: three elements in the augmentation
ideal always multiply to zero, and the remaining cases use the unit.
The pairing is the Frobenius pairing $\tau(xy)$ with $\tau(v)=1$.
It is nondegenerate, graded symmetric, and of degree $-3$.
The only nonzero ternary cyclic tensor is the value on four copies of
$a$.  Its rotation sign in \eqref{cyquant:eq:cyclic-sign} is
$(-1)^{3+3}=1$.

At arity four the only nonunit input that contributes is $(a,a,a,a)$.
Equation \eqref{cyquant:eq:stasheff-sign} becomes
\[
 -m_2(m_3(a,a,a),a)+m_2(a,m_3(a,a,a))
       =-\lambda v+\lambda v=0.
\]
The inner binary terms vanish because $a^2=0$.  With one unit the
nonzero terms cancel in pairs; with more units there is no contribution.
At arity five every composite of two ternary products vanishes, because
the inner output is $b$.  No higher arity can occur in a composite of
the available operations.  This proves all the identities.

If $f_2$ has degree $-1$, then $f_2(a,a)$ is a scalar multiple of $a$.
Every term of its Hochschild coboundary on $(a,a,a)$ is consequently
zero: the inner or outer binary product contains $a^2$.  Thus the value
$m_3(a,a,a)=\lambda b$ cannot be removed by a quadratic change of
coordinates.  An invertible graded linear change only rescales the
one-dimensional $a$ and $b$ spaces.  This proves nonformality.
\end{proof}

For the degree-one element $x=ta$, the cyclic potential is
\[
 \Phi(x)=\sum_{n\geq2}\frac1{n+1}
              \langle m_n(x,\ldots,x),x\rangle
         =\frac\lambda4t^4.
\]
The quartic term records the ternary operation.  It does not force the
binary product to vanish: cyclicity with the unit gives
$\langle m_2(a,b),1\rangle=\langle a,b\rangle=1$.
This finite example makes no assertion of homological smoothness or of
a geometric realization by a compact Calabi--Yau category.  Those are
additional properties of a specified model.

\subsection{Compact conifold objects and their ternary product}

The conifold gives a geometric example in which the cubic relations
determine a nonzero ternary Ext product.  Put
\[
 C=\mathbb C[X,Y,U,V]/(XY-UV),\qquad I=(X,U),\qquad
 A=\operatorname{End}_C(C\oplus I).
\]
Use arrows $a,c:1\to0$ and $b,d:0\to1$, with composition written
right to left, and the potential $\Phi=abcd-adcb$.  The arrow matrices
on $C\oplus I$ are
\[
 a=\begin{pmatrix}0&1\\0&0\end{pmatrix},\quad
 b=\begin{pmatrix}0&0\\X&0\end{pmatrix},\quad
 c=\begin{pmatrix}0&V/X\\0&0\end{pmatrix},\quad
 d=\begin{pmatrix}0&0\\U&0\end{pmatrix}.
\]
Here $V/X$ is a regular homomorphism $I\to C$, since it sends
$X$ to $V$ and $U$ to $Y$.  These matrices identify $A$ with the
Jacobi algebra of $\Phi$.  One way to verify the presentation is to
write $z=a+b$, $p=c+d$, $g=e_0-e_1$, with $e_i$ the vertex
idempotents, and
$\lambda=(cb-ad)+(bc-da)$.  The cyclic derivatives say that
$\lambda$ is central and give
\[
 g^2=1,\quad gz=-zg,\quad gp=-pg,\quad pz-zp=\lambda g.
\]
The normal monomials $z^ip^jg^\varepsilon$ form a
$\mathbb C[\lambda]$-basis: the only mixed ambiguity $gpz$ has the
same reduction $zpg+\lambda$ in both orders.  The central elements
$X=z^2$, $Y=p^2$, $W=zp+\lambda g/2$,
$U=W-\lambda/2$, $V=W+\lambda/2$ satisfy $XY=UV$.
The matrix map is surjective and becomes an isomorphism over
$\operatorname{Frac}(C)$ on the basis $1,z,g,zg$.
Filter the source by the number of $z,p$ factors, giving $\lambda$
weight zero.  Its associated graded is
$\mathbb C[\lambda,z,p]\rtimes\mu_2$; the central graded ring
$\mathbb C[\lambda,X,Y,W]/(XY-W^2)$ embeds by $X=z^2$, $Y=p^2$,
$W=zp$ and acts without torsion.  The filtered source is therefore
torsion-free over $C$, so the kernel is zero.

Let $S_i$ be the right vertex simples and $P_i=e_iA$.  Give the arrows
weight one; in $P_i(-w)$ the free generator has internal weight $w$.
The compact projective resolutions are
\begin{equation}\label{cyquant:eq:conifold-resolutions}
 \begin{aligned}
 0\to P_0(-4)&\xrightarrow{\binom bd}P_1(-3)^2
 \xrightarrow{\left(\begin{smallmatrix}-dc&bc\\da&-ba\end{smallmatrix}\right)}
 P_1(-1)^2\xrightarrow{(a\ c)}P_0\to S_0\to0,\\
 0\to P_1(-4)&\xrightarrow{\binom ac}P_0(-3)^2
 \xrightarrow{\left(\begin{smallmatrix}cd&-ad\\-cb&ab\end{smallmatrix}\right)}
 P_0(-1)^2\xrightarrow{(b\ d)}P_1\to S_1\to0.
 \end{aligned}
\end{equation}
The composites vanish by the four cyclic derivatives.  For exactness,
the first middle matrix in the fractional-ideal model is
$\left(\begin{smallmatrix}-Y&V\\U&-X\end{smallmatrix}\right)$.
Its kernel is $(Xh,Uh)$, with $h$ in $C$ or $I^\vee$ according to
the relevant component.  At the next place the required colon identities
are
\[
 \{n\in C:Vn\in XC\}=I,\qquad
 \{n\in I:Vn\in XI\}=I^2.
\]
The first follows by reducing modulo $X$, where
$C/(X)=\mathbb C[Y,U,V]/(UV)$.  For the second, write
$n=X\alpha+U\beta$ and reduce the condition
$V\alpha+Y\beta\in I$ modulo $I$.  In $\mathbb C[Y,V]$ it gives
$\bar\alpha=Yh$, $\bar\beta=-Vh$, whence $n\in I^2$ by $XY=UV$.
These identities prove exactness of the first resolution; exchange the
vertices for the second.  All entries have positive weight, so the
resolutions are minimal.  They give self-Ext dimensions $(1,0,0,1)$
and off-diagonal dimensions $(0,2,2,0)$.

\begin{proposition}\label{cyquant:prop:conifold-ternary}
A minimal model of
$\operatorname{RHom}_A(S_0\oplus S_1,S_0\oplus S_1)$ has only
$m_2,m_3$ as nonunit operations.  Its ternary product is nonzero and
cannot be removed by a change of minimal model.
\end{proposition}

\begin{proof}
Adjoin opposite arrows $\alpha^*$ of degree $-1$ and loops $\tau_i$
of degree $-2$ to the path algebra, with
\[
 d\alpha^*=\partial_\alpha\Phi,\qquad
 d\tau_i=e_i\sum_\alpha[\alpha,\alpha^*]e_i.
\]
The cyclic derivative identity proves $d^2=0$.  The generating weights
are $1,3,4$.  The one-sided twisted complexes of the dual generating
coalgebra are precisely \eqref{cyquant:eq:conifold-resolutions}.
The universal twisted complex over its cobar resolves the vertex ring;
the positive internal weights make the bar contraction and its
perturbation finite at each weight.  Tensoring it with $A$ gives the
exact complexes above.  If the cone of the cobar map to $A$ had
cohomology, its least internal weight would survive tensoring with the
vertex ring, a contradiction.  Thus the cobar map is a quasi-isomorphism.

The Ext weights are $0,1,3,4$ in cohomological degrees $0,1,2,3$.
Preservation of weight and the rule $|m_n|=2-n$ permit only the
binary products on degree pairs $(1,2),(2,1)$ and the ternary product
on $(1,1,1)$.  In suspended dual coordinates the relation
$d(d^*)=abc-cba$ gives
\[
 b_3(sa^\vee,sb^\vee,sc^\vee)=s(d^*)^\vee,
 \qquad
 b_3(sc^\vee,sb^\vee,sa^\vee)=-s(d^*)^\vee.
\]
One can see the first value directly on the resolutions.  Choose closed
degree-one lifts $A_1,B_1,C_1$ of the three arrow classes, with components
from homological degree $n$ to $n-1$ as follows:
\[
 \begin{array}{c|ccc}
 &n=1&n=2&n=3\\\hline
 A_1&(1,0)&\left(\begin{smallmatrix}0&-c\\c&0\end{smallmatrix}\right)&(1,0)^t\\
 B_1&(1,0)&\left(\begin{smallmatrix}0&d\\-d&0\end{smallmatrix}\right)&(1,0)^t\\
 C_1&(0,1)&\left(\begin{smallmatrix}0&a\\-a&0\end{smallmatrix}\right)&(0,1)^t
 \end{array}
\]
Take homotopies whose only nonzero components are
$H_{AB,2}=\left(\begin{smallmatrix}0&0\\0&1\end{smallmatrix}\right)$
and $H_{BC,2}=\left(\begin{smallmatrix}0&1\\0&0\end{smallmatrix}\right)$.
For $\delta f=d_{\mathrm{target}}f-(-1)^{|f|}fd_{\mathrm{source}}$,
matrix multiplication gives $\delta H_{AB}=A_1B_1$ and
$\delta H_{BC}=B_1C_1$.  The closed representative
$A_1H_{BC}+H_{AB}C_1$ projects to the row $(0,1)$ in degree two.
All products of two degree-one Ext classes vanish by weight, so the
triple Massey value has zero indeterminacy.  It cannot be erased.
\end{proof}

Under the conifold tilting equivalence these objects are
$\mathcal O_E$ and $\mathcal O_E(-1)[1]$, where $E$ is the exceptional curve.
Their degree-three Serre pairing pairs arrows with relations and units
with top classes.  Its compact support is essential: it supplies no
finite trace on all of $A$.  The ternary calculation retains that
pairing and the quartic potential without identifying them with a
quantized field action.

\section{Deforming a specified action}
\label{cyquant:sec:formal-actions}

The equations for a formal action deformation belong to a convolution
complex over a characteristic-zero field $k$.  Fix an augmented dg
algebra $A$, a complex $V$, and a unital
dg action $\rho:A\to E=\operatorname{End}(V)$.  Put
$\bar A=\ker(A\to k)$ and give $s\bar A$ the shift $|sa|=|a|-1$.
On the reduced tensor coalgebra
$B(A)=\bigoplus_{r\geq1}(s\bar A)^{\otimes r}$, use the bar components
\[
 b_1(sa)=-s(da),\qquad b_2(sa\otimes sb)=(-1)^{|a|}s(ab).
\]
Its coproduct is deconcatenation.  The completed convolution algebra is
\[
 \mathcal L=\prod_{r\geq1}
       \operatorname{Hom}((s\bar A)^{\otimes r},E).
\]
For homogeneous $f$, its differential is
$d_{\mathcal L}f=d_Ef-(-1)^{|f|}fb$; its product is composition in $E$
after the coproduct, with the tensor Koszul sign.  The commutator makes
$\mathcal L$ a dg Lie algebra.

The degree-one map $\tau(sa)=\rho(a)$, zero on longer tensors, satisfies
$d_{\mathcal L}\tau+\tau\star\tau=0$.  At tensor length two the two
terms are $(-1)^{|a|}\rho(ab)$ and
$(-1)^{|a|-1}\rho(a)\rho(b)$, so multiplicativity gives their
cancellation.  Set
\[
 D=d_{\mathcal L}+[\tau,-].
\]
A deformation of this action, allowing higher action components, is a
degree-one series $a(t)=\sum_{n\geq1}t^na_n\in t\mathcal L[[t]]$
satisfying
\begin{equation}\label{cyquant:eq:action-mc}
 Da(t)+\tfrac12[a(t),a(t)]=0.
\end{equation}
The $t$-adic completion makes every coefficient a finite sum; each
$a_n$ can itself have arbitrary tensor arity.

\begin{proposition}\label{cyquant:prop:action-obstructions}
Suppose \eqref{cyquant:eq:action-mc} holds through order $n-1$.  Put
\[
 R_n=-\frac12\sum_{i=1}^{n-1}[a_i,a_{n-i}].
\]
Then $DR_n=0$.  An extension to order $n$ exists exactly when
$[R_n]=0$ in $H^2(\mathcal L,D)$.  Its choices form an affine space
over $Z^1(\mathcal L,D)$, and modulo order-$n$ gauge transformations
they form an affine space over $H^1(\mathcal L,D)$.
\end{proposition}

\begin{proof}
The coefficient of $t^n$ in the curvature of $\tau+a(t)$ is
$Da_n-R_n$.  The Bianchi identity says that its twisted differential is
zero.  Since the lower curvature coefficients vanish, the coefficient
of $t^n$ in that identity is $DR_n=0$.  Solving $Da_n=R_n$ gives the
obstruction statement.  Differences of solutions are closed degree-one
elements.  A gauge parameter $t^n\xi$, $|\xi|=0$, changes $a_n$ by
$D\xi$ without changing the lower coefficients, proving the last claim.
\end{proof}

Suppose the twisted complex has a contraction onto its cohomology,
$Dh+hD=1-ip$, with $|h|=-1$.  If $pR_n=0$, the choice $a_n=hR_n$
solves the order-$n$ equation.  Through order six this gives
\begin{equation}\label{cyquant:eq:six-orders}
 \begin{aligned}
 a_2&=-\tfrac12h[a_1,a_1],\\
 a_3&=-h[a_1,a_2],\\
 a_4&=-h\bigl([a_1,a_3]+\tfrac12[a_2,a_2]\bigr),\\
 a_5&=-h\bigl([a_1,a_4]+[a_2,a_3]\bigr),\\
 a_6&=-h\bigl([a_1,a_5]+[a_2,a_4]+\tfrac12[a_3,a_3]\bigr).
 \end{aligned}
\end{equation}
Every line assumes the preceding obstruction projections vanish.  The
factor $1/2$ on the diagonal terms comes from the symmetry of the bracket
on degree-one elements; off-diagonal terms occur twice.

If $a_1=u\beta$ and the contraction is linear over $k[u]$, this recursion
gives $a_n\in u^n\mathcal L$.  The first terms are
\[
 a_2=-\tfrac{u^2}{2}h[\beta,\beta],\qquad
 a_3=\tfrac{u^3}{2}h[\beta,h[\beta,\beta]].
\]
The divisibility is a consequence of the recursion.  It does not specify
the tensor arity of $a_n$, and it need not be the exact order of
vanishing if a coefficient is zero.

Two degree-zero algebras separate extension from obstruction.  For
$A=k[x]$ acting on $k$ by augmentation, the strict deformed action
$x\mapsto tu$ gives
\[
 a_n(sx^r)=\begin{cases}u^n&r=n,\\0&r\ne n,\end{cases}
\]
with no higher tensor-length component.  All orders exist because the
polynomial algebra has no relation to violate.  For
$A=k[\epsilon]/(\epsilon^2)$, the same first-order value
$a_1(s\epsilon)=u$ has a nonzero order-two obstruction: the required
strict relation would give $t^2u^2=0$.  In the convolution complex the
relevant differential is zero and the residual value on
$s\epsilon\otimes s\epsilon$ is nonzero.  No higher action component
with values in $\operatorname{End}(k)$ can remove it.  Thus an
integer sequence attached to another generating problem cannot count
the extensions without an actual identification of complexes.

\section{Compact supports and holomorphic moments}
\label{cyquant:sec:moments}

The support condition changes the coefficient space already in the
linear Dolbeault problem.  On $\mathbb C^3$, write
$z_j=x_j+iy_j$, $dV=\prod_j(dx_j\wedge dy_j)$, and
$\Omega=dz_1\wedge dz_2\wedge dz_3$.  Then
\[
 \Omega\wedge d\bar z_1\wedge d\bar z_2\wedge d\bar z_3=-8i\,dV.
\]
Choose a smooth radial bump $\rho(|z|^2)$ supported inside the unit disk
with $\int_{\mathbb C}\rho\,dA=1$.  For a multi-index
$\beta\in\mathbb N^3$, define compactly supported forms
\begin{equation}\label{cyquant:eq:moment-forms}
 \alpha_0=\frac1{-8i}\prod_j\rho(|z_j|^2)\,
               d\bar z_1\wedge d\bar z_2\wedge d\bar z_3,
 \qquad
 \alpha_\beta=\frac{(-1)^{|\beta|}}{\beta!}
                         \partial_z^\beta\alpha_0.
\end{equation}
They are $\bar\partial$-closed because they have top antiholomorphic
degree.

\begin{proposition}\label{cyquant:prop:moments}
For every entire holomorphic function $f$,
\[
 \int_{\mathbb C^3}f\Omega\wedge\alpha_\beta
       =\frac{\partial_z^\beta f(0)}{\beta!}.
\]
The classes $[\alpha_\beta]$ are linearly independent in
$H_c^{0,3}(\mathbb C^3)$.  For every $N$, the moment map to the dual of
polynomials of total degree at most $N$ is split surjective and has rank
$\binom{N+3}{3}$.
\end{proposition}

\begin{proof}
Angular averaging in each coordinate gives the formula for $\beta=0$.
Repeated integration by parts gives the general formula; compact support
kills every boundary term.  If $\alpha=\bar\partial\eta$ with $\eta$
compactly supported, Stokes gives $\int f\Omega\wedge\alpha=0$.
Pairing a finite relation among the $[\alpha_\beta]$ with $z^\gamma$
therefore extracts its $\gamma$-coefficient.  The same monomial pairing
gives the asserted finite-dimensional splitting.
\end{proof}

There is a useful one-dimensional quotient, but the cohomology itself
is infinite-dimensional.  To identify that quotient, put
$F(t)=-\int_t^1\rho(s)\,ds$, extended by zero for $t\geq1$.
The function $F(|z|^2)$ is smooth and compactly supported, and
$\partial_{\bar z}F(|z|^2)=z\rho(|z|^2)$.  Replacing the $j$th bump
factor in $\alpha_0$ by $F$, omitting $d\bar z_j$, and using the sign
$(-1)^{j-1}$ gives a compactly supported primitive of $z_j\alpha_0$.
Differentiation and the Leibniz rule consequently give
\begin{equation}\label{cyquant:eq:moment-weyl}
 \begin{aligned}
 z_j[\alpha_\beta]&=
 \begin{cases}[\alpha_{\beta-e_j}]&\beta_j>0,\\0&\beta_j=0,\end{cases}\\
 \partial_{z_j}[\alpha_\beta]&=-(\beta_j+1)[\alpha_{\beta+e_j}].
 \end{aligned}
\end{equation}
These operators satisfy $[\partial_{z_j},z_k]=\delta_{jk}$.
The span of the moment classes is the polynomial delta module for the
Weyl algebra.  Its cyclic generator $[\alpha_0]$ is killed by every $z_j$.
Its quotient by $\sum_j\partial_{z_j}$ of that span is the single line
generated by $[\alpha_0]$.  This quotient retains only the zeroth moment;
it is a different coefficient object from $H_c^{0,3}(\mathbb C^3)$.
The compact-support moment mechanism is also visible in the Cauchy
formula of \cite[Ch.~I, \S3]{cyquant:Demailly}.

Gauge coefficients can introduce further dimensions.  If
$\mathfrak g=\bigoplus_{i=1}^r\mathfrak g_i$ is complex semisimple with
simple summands, every invariant symmetric form has zero cross terms:
for $x=[a,b]\in\mathfrak g_i$ and $y\in\mathfrak g_j$, $i\ne j$,
invariance gives $B(x,y)=B(a,[b,y])=0$.  Each simple summand contributes
one line of forms.  A finite-dimensional local deformation claim must
therefore specify both its gauge restriction and its support or
translation quotient, together with the map from the field-theoretic
deformation complex.

\section{A nodal BRST field algebra}
\label{cyquant:sec:nodal}

Let $u,m$ be target coordinates on $\mathbb A^2$, and let $z$ denote
the independent field-insertion coordinate.  Start with even fields
$u,m,\beta_u,\beta_m$ and odd fields $\eta,\chi,b_\eta,b_\chi$.
Their elementary singular products are
\begin{equation}\label{cyquant:eq:nodal-oscillators}
 \beta_u(z)u(w)\sim\frac1{z-w},\quad
 \beta_m(z)m(w)\sim\frac1{z-w},\quad
 b_\eta(z)\eta(w)\sim\frac1{z-w},\quad
 b_\chi(z)\chi(w)\sim\frac1{z-w}.
\end{equation}
Reversing an even pair changes the sign; reversing an odd pair keeps
it.  All other elementary products are regular.

The coordinates have conformal weight zero, their momenta weight one.
The cochain degrees of $\eta,\chi$ are $-1$, those of
$b_\eta,b_\chi$ are $1$, and all even fields have degree zero.
For each coordinate and its momentum use oscillator expansions
$x(z)=\sum_nx[n]z^{-n}$ and
$\beta_x(z)=\sum_n\beta_x[n]z^{-n-1}$.
The vacuum representation has zero-mode space
$\mathbb C[u,m]\otimes\Lambda(\eta,\chi)$; coordinate zero modes
multiply and momentum zero modes differentiate.  Positive modes kill
this space.  Negative modes freely generate polynomial or exterior
factors according to parity.  Normal ordering and Wick contraction
define the vertex superalgebra $V$, including every translated state
\cite[\S\S2.1, 3.2]{cyquant:MalikovSchechtman}.

\begin{proposition}\label{cyquant:prop:nodal-differential}
The charge field and its residue are
\[
 q(z)={:}u(z)m(z)b_\eta(z){:},\qquad
 Q=\operatorname{Res}_{z=0}q(z).
\]
The operator $Q$ is a square-zero, weight-preserving vertex derivation of degree one.
It kills the vacuum, commutes with translation, and acts by
\[
 Q\eta=um,\qquad Q\beta_u=-{:}mb_\eta{:},\qquad
 Q\beta_m=-{:}ub_\eta{:},
\]
with zero value on $u,m,\chi,b_\eta,b_\chi$.
Its weight-zero cohomology is
$\mathbb C[u,m]/(um)\otimes\Lambda(\chi)$, and its vacuum is nonzero
in the full cohomology.
\end{proposition}

\begin{proof}
A zero mode is a derivation of every vertex product, kills the vacuum,
and commutes with translation.  The charge has no self-contractions,
so its self-OPE is regular and $2Q^2=[Q,Q]=0$.  Its weight is one,
making its residue weight zero.  Wick contraction gives the displayed
values, with both minus signs coming from the reversed even pairs.
In weight zero the complex is
$(\mathbb C[u,m]\otimes\Lambda(\eta,\chi),um\partial_\eta)$.
Multiplication by $um$ is injective in the polynomial ring, so its
cohomology is the stated quotient and the unit survives.  Weight
preservation prevents a higher-weight primitive from killing it.
\end{proof}

Every translated relation is retained:
\[
 Q(T^n\eta)=\sum_{a=0}^n\binom na{:}(T^au)(T^{n-a}m){:}.
\]
The momentum zero mode $\beta_{u,(0)}$ is not closed.  The construction
therefore makes normal differentiation part of a dg operator action;
it does not turn it into a closed normal connection on the unchanged
nodal quotient.

Define closed currents
\[
 J_u={:}u\beta_u{:}+{:}\eta b_\eta{:},\qquad
 J_m={:}m\beta_m{:}+{:}\eta b_\eta{:},\qquad
 J_\chi={:}\chi b_\chi{:}.
\]
Their pairing matrix, in this order, is
\begin{equation}\label{cyquant:eq:nodal-current-pairing}
 J_i(z)J_j(w)\sim\frac{B_{ij}}{(z-w)^2},\qquad
 B=\begin{pmatrix}0&1&0\\1&0&0\\0&0&1\end{pmatrix}.
\end{equation}
Indeed the two terms of $QJ_u$, and of $QJ_m$, cancel.  A bosonic
number current has double contraction $-(z-w)^{-2}$ and an odd number
current has $+(z-w)^{-2}$; the single contractions cancel.  The two
mixed currents share only the odd pair.  This proves the formula.
It embeds the rank-three Heisenberg vacuum algebra in $H(V,Q)$:
the nondegenerate pairing lets positive modes differentiate a highest
negative-mode monomial to a nonzero scalar, so the Heisenberg vacuum
module is simple.  Its map to cohomology preserves the nonzero vacuum
and is therefore injective.

\subsection{A rational dressing with a finite pole bound}

Work over $R_N=\mathbb C[a,b,\varepsilon]/(\varepsilon^N)$ and expand
first in $\varepsilon$.  Set
\[
 K(s)=\frac{s(s+\varepsilon(a+b))}
              {(s+\varepsilon a)(s+\varepsilon b)},\qquad
 r(s)=K(s)/K(-s).
\]
Write $J_u(z)=\sum_nJ_{u,n}z^{-n-1}$.
For $g(z)\in1+\varepsilon R_N[z^{-1}]$, define
\begin{equation}\label{cyquant:eq:dressing}
 A_g(z)=\sum_{n\geq0}\frac{(-1)^n}{n!}
        \partial_z^n\log g(z)\,J_{u,n},\qquad U_g(z)=e^{A_g(z)}.
\end{equation}
The logarithm and exponential terminate in $\varepsilon$.  On a state
of weight $L$, the current modes with $n>L$ vanish, and the other
nonnegative modes do not increase weight.  Thus the formula defines an
operator on every state.

The modes $J_{u,n}$ commute with one another and with $Q$.  Their action
$[J_{u,n},Y(u,w)]=w^nY(u,w)$ gives
\[
 U_g(z)Y(u,w)U_g(z)^{-1}=g(z-w)Y(u,w),\qquad |z|>|w|.
\]
The coefficient relation $a_n'=-(n+1)a_{n+1}$ in
\eqref{cyquant:eq:dressing}, together with
$[T,J_{u,n}]=-nJ_{u,n-1}$, gives $[T,U_g]=\partial_zU_g$.
Every $U_g$ fixes the vacuum.  Consequently
$C(z)=Y(u,z)U_K(z)$ and $H(p)=U_r(p)$ obey
\begin{equation}\label{cyquant:eq:nodal-exchange}
 \begin{aligned}
 C(z)C(w)&=r(z-w)C(w)C(z),\\
 H(p)C(z)H(p)^{-1}&=r(p-z)C(z),\\
 [Q,C(z)]&=[Q,H(p)]=0,\qquad C(z)\mathbf1=e^{zT}u.
 \end{aligned}
\end{equation}
The equalities compare rational continuations and their specified
Laurent expansions.  For any finite number of insertions,
\[
 C(z_1)\cdots C(z_r)=
 \prod_{i<j}K(z_i-z_j)\,
 Y(u,z_1)\cdots Y(u,z_r)\prod_iU_K(z_i).
\]
At order $\varepsilon^j$ the pair-kernel product has total diagonal
pole order at most $j$; derivatives of total order $M$ add at most
$M$.  A fixed oscillator state contributes its own finite pole bound.
These bounds give compatible operators for all $N$.

For example, modulo $\varepsilon^4$,
\[
 K(s)=1-\frac{ab\varepsilon^2}{s^2}
       +\frac{ab(a+b)\varepsilon^3}{s^3},\qquad
 r(s)=1+\frac{2ab(a+b)\varepsilon^3}{s^3}.
\]
This constructs the displayed operator family.  A new state-field map
on every state with these assigned operators has not been defined by
the dressing alone.  Its comparison with an oriented Hall operation
requires the corresponding morphism of coefficients and residues.

\subsection{Induction from differential modules}

Put $A_0=\mathbb C[u,m]$ and let $\mathcal D_{A_0}$ be its polynomial
differential-operator algebra.  The super differential-operator algebra
of $P=A_0\otimes\Lambda(\eta,\chi)$ has differential
$[um\partial_\eta,-]$.  In particular,
\[
 [um\partial_\eta,\partial_u]=-m\partial_\eta,
 \qquad [um\partial_\eta,\partial_m]=-u\partial_\eta.
\]
For a complex $(M,d_M)$ of left $\mathcal D_{A_0}$-modules, set
\[
 L(M)=M\otimes\Lambda(\eta,\chi),\qquad
 D_M=d_M+um\partial_\eta
\]
with the graded tensor sign.  Its square is zero because $d_M$ is
$\mathcal D_{A_0}$-linear.
Let $\widehat{\mathfrak o}$ be the oscillator Lie superalgebra of
\eqref{cyquant:eq:nodal-oscillators}, with central element acting by
one.  Give $L(M)$ its indicated zero-mode action and let strictly
positive modes act by zero.  Oscillator induction gives
\begin{equation}\label{cyquant:eq:nodal-induction}
 \mathbb M(M)=U(\widehat{\mathfrak o}_{<0})\otimes L(M),
 \qquad d_{\mathbb M}=d_M+q_{(0)}.
\end{equation}
PBW supplies polynomial factors for even negative modes and exterior
factors for odd ones.  Normal-ordered state fields act with finitely
many terms on each vector.  On the top space, $q_{(0)}$ is
$um\partial_\eta$; on the full induced module its square vanishes by
the same regular self-OPE.  Thus \eqref{cyquant:eq:nodal-induction}
defines a dg $V$-module.

A $\mathcal D_{A_0}$-linear chain map extends by the identity on all
oscillator factors.  It commutes with every state-field mode, with the
dressings, and with principal-part residues and their finite normal
derivatives.  Induction preserves short exact sequences of underlying
complexes because these factors are free vector spaces.  Positive
oscillator modes act locally nilpotently by differentiating finite
monomials.  A spectator connection commuting with the ambient action
and $d_M$ and annihilating $um$ extends through the same formulas.

\subsection{The constant critical coefficient}

For $f=um$, let
\[
 N=A_0[f^{-1}]/A_0,\quad N_u=A_0[u^{-1}]/A_0,\quad
 N_m=A_0[m^{-1}]/A_0.
\]
Their quotient sequence
\[
 0\longrightarrow N_u\oplus N_m\longrightarrow N
              \longrightarrow\delta_0\longrightarrow0
\]
has $\delta_0$ spanned by $u^{-r}m^{-s}$, $r,s\geq1$.
It is nonsplit even as an $A_0$-module sequence.  Indeed,
$\ker(u:N\to N)$ is spanned by $u^{-1}m^b$, $b\geq0$, while
$\ker(m:N\to N)$ is spanned by $u^am^{-1}$, $a\geq0$.
Their intersection is zero; the nonzero class $u^{-1}m^{-1}$ in
$\delta_0$ is killed by both.  Induction carries the sequence and
its actual maps to an exact sequence of dg field modules.  This
exactness makes no assertion that every derived extension remains
nonsplit.

The constant critical coefficient is the holonomic dual
$P_Z=\mathbb D_{\mathcal D}N$.  With the normalization
$\operatorname{DR}(\mathcal O_{\mathbb A^2})=\mathbb C_{\mathbb A^2}[2]$,
regular-holonomic duality gives
$\operatorname{DR}(P_Z)=i_{Z,*}\mathbb C_Z[1]$ for $Z=\{um=0\}$.
There is an explicit presentation:
\begin{equation}\label{cyquant:eq:critical-module}
 P_Z=\ker\left(
 J=\mathcal D_{A_0}/(\mathcal D_{A_0}u\partial_u+
                    \mathcal D_{A_0}m\partial_m)
       \xrightarrow{\ g\mapsto1\ }A_0\right).
\end{equation}
Here $g$ is the class of one in $J$.
To see the presentation, $A_0[f^{-1}]$ is cyclic on $1/(um)$ with commuting
relations $u\partial_u+1$ and $m\partial_m+1$.  Their principal
symbols form a regular sequence.  Formal transposition in its Koszul
resolution sends them to $-u\partial_u$ and $-m\partial_m$.
Dualizing the localization sequence proves
\eqref{cyquant:eq:critical-module}.

A basis of its kernel is
\[
 A_{r,b}=\partial_u^rm^bg\ (r\geq1,b\geq0),\quad
 B_{a,s}=u^a\partial_m^sg\ (a\geq0,s\geq1),\quad
 C_{r,s}=\partial_u^r\partial_m^sg\ (r,s\geq1).
\]
The polynomial--polynomial sector maps isomorphically to $A_0$;
these other three sectors give the kernel.  The equality
$\partial_mA_{r,0}=C_{r,1}$ retains the attachment of a branch to
the crossing.
Commuting the coordinates through the derivatives gives
\begin{equation}\label{cyquant:eq:critical-kernel}
 \begin{aligned}
 fA_{r,b}&=-(r-1)A_{r-1,b+1},\\
 fB_{a,s}&=-(s-1)B_{a+1,s-1},\\
 fC_{r,s}&=(r-1)(s-1)C_{r-1,s-1},
 \end{aligned}
\end{equation}
with zero whenever a coefficient vanishes.
The cokernel has basis $A_{r,0},B_{0,s}$, and the kernel has basis
$A_{1,b},B_{a,1},C_{1,s},C_{r,1}$, counting $C_{1,1}$ once.
Every nonzero image in each sector is a distinct basis vector, which
proves these assertions.  Hence
\[
 H(L(P_Z))=
  \bigl(\operatorname{coker}f\oplus\eta\operatorname{ker}f\bigr)
       \otimes\Lambda(\chi).
\]
The first two summands have degrees zero and $-1$.  On its degree-zero
top cohomology, $(J_{u,0},J_{m,0})$ has weights $(-r,0)$ on $A_{r,0}$
and $(0,-s)$ on $B_{0,s}$.  This is a concrete coefficient module,
distinct from the vacuum representation.

For comparison, $f:N\to N$ is surjective with kernel
$A_0/(f)$, via $a\mapsto[a/f]$.  Thus
$H^{-1}(L(N))=A_0/(f)$, $H^{-2}(L(N))=\chi A_0/(f)$, and
$H^0(L(N))=0$.  The field-module construction preserves the actual
dg differential-operator action, including its nonclosed normal
momenta, rather than identifying these two cohomologies.

Finally a moving spectator coefficient can be added with generator
$\delta_\lambda$ satisfying
\[
 (w-\lambda)\delta_\lambda=0,\qquad
 (\partial_w+\partial_\lambda)\delta_\lambda=0.
\]
The two operators commute.  Tensoring it with the nodal coefficient
retains this common connection under induction, since $um$ is
independent of $w,\lambda$.  If
$\lambda=(a_{\mathrm{fr}}+\varepsilon y)/\beta$ with $\beta$ invertible,
the common derivative is $\partial_w+\beta\partial_{a_{\mathrm{fr}}}$.
Its identification with a framed Hall direct image still requires the
oriented correspondence map.

\section{The interacting curved dual of the Miura algebra}
\label{cyquant:sec:wfour}

The next example has noncentral ghost brackets.  Work over
$k=\mathbb C(\alpha)$ on $X=\mathbb A^1_k$, with translation $T=\partial_z$.
Let four even Heisenberg fields satisfy
\[
 \sum_{i=1}^4a_i=0,\qquad
 a_i(z)a_j(w)\sim\frac{\delta_{ij}-1/4}{(z-w)^2}.
\]
Define $U_s$ by the right-nested, normally ordered Miura operator
\begin{equation}\label{cyquant:eq:miura}
 {:}(\alpha T+a_1)(\alpha T+a_2)(\alpha T+a_3)(\alpha T+a_4){:}
   =\sum_{s=0}^4U_s(\alpha T)^{4-s},
 \qquad U_0=\mathbf1,\quad U_1=0,
\end{equation}
moving every $T$ to the right.  Let $V$ be the generated vertex
subalgebra.  Its stress field is $L=-U_2$ and its central charge is
$c=3-60\alpha^2$.  Relative to the negative Heisenberg contraction in
\cite{cyquant:Zhu}, the convention change is
$a_i^Z=ia_i$, $\alpha_Z=i\alpha$, $U_s^Z=i^sU_s$.

All coefficients can be computed with one Wick rule:
\begin{equation}\label{cyquant:eq:miura-wick}
 (T^ra_i)(z)(T^sa_j)(w)\sim
  (\delta_{ij}-1/4)(-1)^r(r+s+1)!(z-w)^{-r-s-2}.
\end{equation}
Sum over partial cross-matchings and Taylor-expand each uncontracted
left factor.  This includes regular terms that become poles after
multiplication by a coefficient.  For instance, if
$\operatorname{PP}Y(U_4,u)U_4=\sum_{p=1}^8A_pu^{-p}$, then
\[
 A_8=\frac{c(5-180\alpha^2+2016\alpha^4)}{32}\mathbf1,
 \qquad A_6=-\frac54(1-12\alpha^2)^2U_2.
\]
The coefficient $A_2$ contains $\frac34(U_3)_{(-1)}U_3$.
It cannot be discarded: the $a_1^4a_2^2$ coefficient detects it.
The odd coefficients are fixed by self-skewsymmetry; for example,
$A_5=\frac12TA_6$ and
$A_3=\frac12TA_4-\frac1{24}T^3A_6$.

\subsection{The complete finite relation datum}

Put $I_4=\{0,2,3,4\}$ and let $N^\circ$ have even frame $e_i$,
with $e_i\mapsto U_i$.  The index $0$ is a formal scalar generator.
The right $\mathcal D$ generators are $N^\circ\otimes\omega$.
Their shifted dual is
\[
 \check N^\circ=s^{-1}(N^\circ)^\vee\otimes\omega^{-1},
 \qquad
 \langle-,-\rangle:(N^\circ\boxtimes N^\circ)
       \otimes(\check N^\circ\boxtimes\check N^\circ)
       \longrightarrow s^{-2}\omega^{-1}_{X^2}.
\]
Here $s^{-1}$ raises degree by one.  Write its odd frame as
$\Gamma_0=S$, $\Gamma_s=\gamma_s$ for $s=2,3,4$;
all four generators have cochain degree one and weight $1-s$.
In the affine coordinate their density types are
$s^{-1}\omega^{s-1}$, including $s^{-1}\omega^{-1}$ at $s=0$.

The following finite matrix specifies both relation lattices.
For $u=z_1-z_2$, define
\begin{equation}\label{cyquant:eq:wfour-dual-frame}
 t_{ij}=\Gamma_i\boxtimes\Gamma_j+
       \sum_{k,l}E_{ij}^{kl}u^{-(k+l-i-j)}
                            \Gamma_k\boxtimes\Gamma_l,
 \qquad t=D(\Gamma\boxtimes\Gamma).
\end{equation}
The independent nonzero coefficients are in
Table~\ref{cyquant:tab:wfour}.  Set all others to zero and use
$E_{ji}^{lk}=(-1)^{k+l-i-j}E_{ij}^{kl}$ for reversed rows.
Let $M=D^{-\mathrm t}$ and $r=M(e\boxtimes e)$.
The $\mathcal O_{X^2}$-span $P^\circ$ of the rows $r$ is the primal
quadratic--linear--scalar relation lattice; the span of $t$ is its
annihilator for the displayed density-valued pairing.

\begingroup\small\renewcommand{\arraystretch}{1.13}
\begin{longtable}{@{}cccl@{}}
\caption{The annihilator matrix; $m=k+l-i-j$.}
\label{cyquant:tab:wfour}\\
\toprule $ij$&$kl$&$m$&$E_{ij}^{kl}$\\\midrule\endfirsthead
\toprule $ij$&$kl$&$m$&$E_{ij}^{kl}$\\\midrule\endhead
\bottomrule\endfoot
$00$&$22$&$4$&$\frac32-30\alpha^2$\\
$00$&$23$&$5$&$-6\alpha(20\alpha^2-1)$\\
$00$&$24$&$6$&$-9\alpha^2(20\alpha^2-1)$\\
$00$&$32$&$5$&$6\alpha(20\alpha^2-1)$\\
$00$&$33$&$6$&$(36\alpha^2-1)(20\alpha^2-1)$\\
$00$&$34$&$7$&$3\alpha(20\alpha^2-1)(24\alpha^2-1)$\\
$00$&$42$&$6$&$-9\alpha^2(20\alpha^2-1)$\\
$00$&$43$&$7$&$-3\alpha(20\alpha^2-1)(24\alpha^2-1)$\\
$00$&$44$&$8$&$(3-60\alpha^2)(2016\alpha^4-180\alpha^2+5)/32$\\
$02$&$22$&$2$&$-1$\\
$02$&$23$&$3$&$-4\alpha$\\
$02$&$24$&$4$&$(1-36\alpha^2)/4$\\
$02$&$33$&$4$&$12\alpha^2-1$\\
$02$&$34$&$5$&$4\alpha(12\alpha^2-1)$\\
$02$&$43$&$5$&$-\alpha(12\alpha^2-1)$\\
$02$&$44$&$6$&$-5(12\alpha^2-1)^2/8$\\
$03$&$23$&$2$&$-2$\\
$03$&$24$&$3$&$-3\alpha$\\
$03$&$32$&$2$&$-1$\\
$03$&$33$&$3$&$2\alpha$\\
$03$&$34$&$4$&$12\alpha^2-1$\\
$03$&$43$&$4$&$-1/2$\\
$03$&$44$&$5$&$-3\alpha(6\alpha^2-1)$\\
$04$&$24$&$2$&$-3$\\
$04$&$33$&$2$&$-2$\\
$04$&$34$&$3$&$6\alpha$\\
$04$&$42$&$2$&$-1$\\
$04$&$43$&$3$&$2\alpha$\\
$04$&$44$&$4$&$(3-6\alpha)(2\alpha+1)/2$\\
$22$&$33$&$2$&$1$\\
$22$&$34$&$3$&$\alpha$\\
$22$&$43$&$3$&$-\alpha$\\
$22$&$44$&$4$&$5/8-9\alpha^2/2$\\
$23$&$34$&$2$&$1/2$\\
$23$&$44$&$3$&$-3\alpha/2$\\
$24$&$44$&$2$&$-1$\\
$33$&$44$&$2$&$3/4$\\
\end{longtable}
\endgroup

\begin{proposition}\label{cyquant:prop:wfour-primal}
The lattice $P^\circ$ is free of rank sixteen, agrees with the full
generator tensor space away from the diagonal, and satisfies the unit
intersection condition.  Its unital vertex quotient is $V$.  The rows
of \eqref{cyquant:eq:wfour-dual-frame} are exactly its annihilator,
and the nine nonunit rows give the annihilator of its projected
quadratic lattice.
\end{proposition}

\begin{proof}
Order pairs by $i+j$.  The matrix $D$ is unipotent triangular; each
correction raises the pair sum by at least two.  Its inverse transpose
has the opposite triangularity.  With unit rows first it has block form
\[
 M=\begin{pmatrix}1&0\\ L_0&A_0\end{pmatrix},\qquad
 D=\begin{pmatrix}1&-L_0^{\mathrm t}A_0^{-\mathrm t}\\0&A_0^{-\mathrm t}\end{pmatrix}.
\]
Both diagonal blocks are unipotent.  A linear combination of primal
rows whose nonunit projection vanishes therefore has zero coefficients
on all nine nonunit rows.  Its seven remaining coefficients are regular.
This is the unit intersection condition.

Put $F_{ij}(u)=Y(U_i,u)U_j$.  Applying the Wick rule
\eqref{cyquant:eq:miura-wick} gives $\operatorname{PP}(MF)=0$.
This is an exact finite polynomial calculation: each contraction lowers
the Heisenberg degree, and only the Taylor terms that can reach a
negative power of $u$ are needed.  In particular the scalar coefficient
of $r_{44}$ includes the poles of lower full products; it is not the
isolated coefficient $A_8$ of the self-OPE.

Give $U_s$ filtration $s$ and $T$ filtration zero.  Every relation
expresses a singular product in smaller filtration, so ordered normal
monomials in the jets $T^rU_s$ span the quotient.  In the Heisenberg
realization the leading symbols are the jets of the elementary
symmetric polynomials $e_2,e_3,e_4$ on $\sum a_i=0$.  Their Jacobian
is nonzero at distinct roots.  At every finite jet order its prolongation
is block triangular with these same invertible diagonal blocks.
The jets are algebraically independent, so the spanning monomials
remain independent in $V$.  This proves the presentation.

Finally $MD^{\mathrm t}=1$.  Pairing regularly against every primal
row is equivalent to having regular coordinates in the dual row
frame.  This proves equality with the annihilator, including its
density factor.  The block formula proves the projected assertion.
\end{proof}

\subsection{A finite recursion and the full PBW algebra}

For every $n\geq0$, impose
\begin{equation}\label{cyquant:eq:wfour-mode-rule}
 (\Gamma_i)_{(n)}\Gamma_j
   =-\sum_{kl}E_{ij}^{kl}(\Gamma_k)_{(n-m)}\Gamma_l,
       \qquad m=k+l-i-j.
\end{equation}
For $r\geq0$ use
$(\Gamma_k)_{(-r-1)}\Gamma_l=(T^r\Gamma_k)_{(-1)}\Gamma_l/r!$;
for remaining nonnegative indices apply the rule again.  The positive
filtration
\[
 \delta(S)=9,\quad\delta(\gamma_2)=5,\quad
 \delta(\gamma_3)=3,\quad\delta(\gamma_4)=1,\quad\delta(T)=0
\]
drops by $2m\geq4$ at each correction.  Thus the recursion terminates.

\begin{proposition}\label{cyquant:prop:wfour-pbw}
The rule \eqref{cyquant:eq:wfour-mode-rule} defines a vertex
superalgebra $B^\circ$ with the ordered odd-jet PBW basis.  The
subalgebra $B$ generated by $\gamma_2,\gamma_3,\gamma_4$ embeds in it.
Every nonnegative product $S_{(n)}B$ and $S_{(n)}S$ lies in $B$.
\end{proposition}

\begin{proof}
We specify the finite generator verification and its extension to all
states.  Reduce words by the supercommutator and iterate formulas
\begin{align}
 [x_{(m)},y_{(n)}\}&=
       \sum_{r\geq0}\binom mr(x_{(r)}y)_{(m+n-r)},
       \label{cyquant:eq:wfour-comm}\\
 (x_{(-p)}v)_{(n)}b&=
 \sum_{r\geq0}\binom{p+r-1}{r}
 \left(x_{(-p-r)}v_{(n+r)}b
       -(-1)^{p+|v|}v_{(n-p-r)}x_{(r)}b\right),
       \label{cyquant:eq:wfour-iterate}
\end{align}
where $x$ is odd and $p\geq1$.  Exchange wrong-ordered creation
pairs with the odd sign and their commutator; replace an equal pair
by half its commutator.  Each correction has lower filtration.
For a differentiated generator,
\[
 (\Gamma_{i,(-p)}\mathbf1)_{(n)}
    =(-1)^{p-1}\binom n{p-1}(\Gamma_i)_{(n-p+1)}.
\]
Binomial coefficients with negative upper entry have their polynomial
meaning.  Translation replaces a creation index $p$ by $p+1$ and
multiplies by $p$.

A degree-$d$ normal word has weight at least the sum of the $d$
smallest entries in $\{1-i+r:i\in I_4,r\geq0\}$.
For degrees two and three these bounds are $-5$ and $-7$.
Thus the generator skew identities have only
$0\leq n\leq6-i-j$ to consider, and generator Jacobi only
$m,n\geq0$, $m+n\leq8-i-j-k$.  Substitution of
Table~\ref{cyquant:tab:wfour} into
\[
 \begin{aligned}
 K_{ij,n}&=(\Gamma_i)_{(n)}\Gamma_j-
  \sum_{r\geq0}\frac{(-1)^{n+r}}{r!}T^r(\Gamma_j)_{(n+r)}\Gamma_i,\\
 J_{ijk;mn}&=(\Gamma_i)_{(m)}((\Gamma_j)_{(n)}\Gamma_k)
 +(\Gamma_j)_{(n)}((\Gamma_i)_{(m)}\Gamma_k)\\
 &\qquad-\sum_{r=0}^{m}\binom mr
       ((\Gamma_i)_{(r)}\Gamma_j)_{(m+n-r)}\Gamma_k
 \end{aligned}
\]
reduces each coefficient to zero over $\mathbb Q[\alpha]$ by
\eqref{cyquant:eq:wfour-comm}--\eqref{cyquant:eq:wfour-iterate}.
There are respectively $41$ and $438$ identities.  The weight bounds
exclude every other generator case; reductions use only the relation
module in lower filtration, so confluence has not been assumed.

Choose tensor-algebra lifts of the resulting brackets and replace
$P_{ij}(\lambda)$ by
$\frac12(P_{ij}(\lambda)+P_{ji}(-\lambda-T))$.
The difference is already a lower-filtration relation.  Induction on
the positive filtration preserves both the relation module and the
Jacobi congruences.  These exactly skew lifts meet the nonlinear PBW
hypotheses of \cite[Theorem~3.9]{cyquant:DeSoleKac}.  That theorem
gives the vertex algebra and its ordered odd-jet basis.  All corrections
in the table use indices $2,3,4$, so the three-generator algebra closes,
its PBW words remain independent, and all asserted scalar-dual products
lie in it.  The commutator and iterate formulas extend this preservation
to every state.
\end{proof}

The algebra $B$ is noncommutative.  If $\chi=\gamma_4$, the recursion
gives
\[
 (\gamma_2)_{(0)}\chi=-\chi_{(-1)}\chi_{(-2)}\mathbf1,
 \qquad
 (\gamma_3)_{(0)}\gamma_3=\tfrac34\chi_{(-1)}\chi_{(-2)}\mathbf1.
\]
The common ordered word is nonzero by PBW.

\subsection{Curvature and the exact coupling}

Set
\[
 d=S_{(0)}|_B,\qquad w=\tfrac12S_{(0)}S,
 \qquad q=\sum_{s=2}^4U_s\otimes\gamma_s.
\]
The zero-mode commutator gives $d^2=w_{(0)}$ on every state.  Put
$A_S=S_{(0)}S$.  Jacobi gives $2S_{(0)}A_S=(A_S)_{(0)}S$;
super-skewsymmetry then gives the explicit Bianchi primitive
\begin{equation}\label{cyquant:eq:wfour-bianchi}
 dw=Tv,\qquad
 v=\frac16\sum_{j=1}^{8}\frac{(-1)^{j+1}}{j!}
              T^{j-1}(S_{(j)}A_S).
\end{equation}
The upper limit follows from the degree-three minimum weight $-7$.
The action of the curvature is nonzero.  For example,
\begin{equation}\label{cyquant:eq:wfour-curvature}
 \begin{aligned}
 w_{(0)}\chi={}&\frac c2\gamma_{2,(-4)}\chi_{(-1)}\chi_{(-2)}\mathbf1
 -2c\alpha\gamma_{3,(-5)}\chi_{(-1)}\chi_{(-2)}\mathbf1\\
 &+3c\alpha^2\chi_{(-1)}\chi_{(-2)}\chi_{(-6)}\mathbf1.
 \end{aligned}
\end{equation}
The three PBW words are distinct and $c\ne0$ in $k$.

\begin{theorem}\label{cyquant:thm:wfour-coupling}
With $e=S+q$, one has $e_{(n)}e=0$ for every $n\geq0$.
Consequently $Q=d+q_{(0)}$ is a square-zero differential on $V\otimes B$.
For any unital vertex superalgebra $C$ concentrated in cohomological
degree zero, homomorphisms $V\to C$ are naturally equivalent to tuples
of even states $(x_2,x_3,x_4)$ such that
\[
 (S+\textstyle\sum_sx_s\otimes\gamma_s)_{(n)}
 (S+\textstyle\sum_sx_s\otimes\gamma_s)=0\quad(n\geq0).
\]
The assertion concerns this generator-linear form of coupling.
\end{theorem}

\begin{proof}
Let $F_{ij}=Y(U_i,u)U_j$ and $G_{ij}=Y(\Gamma_i,u)\Gamma_j$.
The primal generators are even, so the tensor OPE has no extra cross
sign.  The dual-frame identity gives
\[
 Y(e,u)e=F^{\mathrm t}G=(MF)^{\mathrm t}(M^{-\mathrm t}G).
\]
Both factors on the right are regular by their defining lattices.
This proves the full chiral identity.  Its zero mode gives
$2e_{(0)}^2=(e_{(0)}e)_{(0)}=0$, and its restriction to
$V\otimes B$ is $Q$.

For the converse, put $x_0=\mathbf1$,
$R=M(Y(x_i,u)x_j)_{ij}$, and $H=M^{-\mathrm t}G$.
The seven unit components of $R$ are regular and
$R_{ji}(u)=e^{uT_C}R_{ij}(-u)$.
The leading dual-filtration symbol of $H_{ij}$ is
$(e^{uT}\xi_i)\xi_j$ for independent odd differential generators
$\xi_i$.

We use their jet independence.  If
$\operatorname{PP}\sum R_{ij}(u)(e^{uT}\xi_i)\xi_j=0$, the coefficient
of $(T\xi_i)\xi_j$ shows that every $R_{ij}$ has at most a simple
pole.  For $i<j$ the coefficient of $\xi_i\xi_j$ and the transported
skew identity give twice its residue; for $i=j$ the same skew identity
makes its residue equal its negative.  All residues vanish in
characteristic zero.  This reasoning also applies to a fixed collection
of unordered pairs.

Apply it to the highest filtration of the coupling equation.  First
$R_{22}$ is regular.  Its entire contribution may then be removed
from principal parts because $H_{22}$ is regular.  Repeat at pair sums
$5,6,7,8$.  Distinct unordered ghost pairs are independent, so all nine
nonunit components become regular.  The tuple therefore satisfies the
complete primal relation lattice and extends uniquely to $V\to C$.
A homomorphism conversely transports the coupling identity.  No
conformal grading on $C$ is required: its trial coupling must satisfy
all nonnegative modes, with locality giving a finite principal part
for each tuple.
\end{proof}

This argument does not identify the dual lattice with the entire kernel
of its chiral operation.  In fact $u^{-1}\chi\boxtimes\chi$ is in
that kernel, since $\chi_{(n)}\chi=0$ for $n\geq0$ and
$\chi_{(-1)}\chi=0$.  Its coefficient in the $44$ relation row is
$u^{-1}$, so it is outside the lattice.  The recognition theorem uses
independent odd jets, which retain more information than this kernel
comparison.

\subsection{Primary coordinates and a translation homotopy}

Set $\Delta_4=22+5c=37-300\alpha^2$ and
\[
 d_0=\frac{480\alpha^4-124\alpha^2+3}{8(300\alpha^2-37)},\qquad
 e_0=-\frac{108\alpha^2-5}{4(300\alpha^2-37)}.
\]
The triangular primary generators are
\[
 W=U_3-\alpha TU_2,\qquad
 H=U_4-\tfrac\alpha2TU_3+d_0T^2U_2+e_0{:}U_2U_2{:}.
\]
The Wick rule gives $L_{(0)}W=TW$, $L_{(1)}W=3W$, and
$L_{(n)}W=0$ for $n\geq2$, with the analogous identities of
weight four for $H$.  The linear system
for the three corrections to $U_4$ has a minor
$-24(300\alpha^2-37)$, proving uniqueness of this ansatz.  The change
is invertible over $\mathbb C[\alpha,\Delta_4^{-1}]$; the Miura
lattices themselves had coefficients in $\mathbb Q[\alpha]$.
The inverse gives the exact coupling
\begin{equation}\label{cyquant:eq:primary-coupling}
 \begin{aligned}
 q={}&-L\gamma_2+W\gamma_3-\alpha(TL)\gamma_3
       +H\gamma_4+\tfrac\alpha2(TW)\gamma_4\\
 &+(d_0-\tfrac{\alpha^2}{2})(T^2L)\gamma_4
          -e_0{:}LL{:}\gamma_4.
 \end{aligned}
\end{equation}
Products between primal and ghost factors denote tensor products.
Integrating derivatives by parts changes this state by an explicit
total derivative.  Such a change preserves its zero mode but need
not preserve its full chiral coupling.  The nonlinear term in
\eqref{cyquant:eq:primary-coupling} is also part of the comparison.

There is nevertheless an exact homotopy for translation.  Define a
degree-$-1$ derivation on $B$ by
\[
 h\gamma_2=-\mathbf1,\qquad h\gamma_3=h\gamma_4=0.
\]
It exists because $\gamma_2\mapsto\gamma_2-\varepsilon\mathbf1$,
with $|\varepsilon|=1$, $\varepsilon^2=0$ and $T\varepsilon=0$, preserves
every generator singular product: their right sides involve only jets
of $\gamma_3,\gamma_4$.  Taking the coefficient of that constant in
the resulting automorphism gives $h$.  It commutes with $T$ and has
$h^2=0$.
In $d\gamma_i$ the only term with undifferentiated $\gamma_2$ is
$-\gamma_{2,(-1)}T\gamma_i$.  The recursion therefore gives
$[d,h]\gamma_i=T\gamma_i$ and hence $[d,h]=T_B$ on every state.
Extend $h$ by zero on $V$.  The exact coupling gives $hq=-U_2=L$,
so
\begin{equation}\label{cyquant:eq:cartan}
 [Q,h]=T_V+T_B.
\end{equation}
For $Q$-cocycles $x,y$ and $n\ne-1$, sesquilinearity gives
\[
 x_{(n)}y
 =Q\left(-\frac1{n+1}(hx)_{(n+1)}y\right).
\]
Thus cohomology has only its $(-1)$-product, which is unital,
graded-commutative, and associative by the vertex identities.
This proves an algebraic translation homotopy.  A realization by a
theory on $\mathbb R\times X$ must additionally compare the full
coupling and the operations on collision strata.  A cooperation on
the curved field dual and its comparison with the completed ordered
bar dual likewise require maps preserving all those operations, the
nonempty vacuum, and the bar unit.

\section{Completed rational ghost modes}
\label{cyquant:sec:ghost-modes}

Return to rational exchange and put
$R_N=\mathbb C[\hbar]/(\hbar^N)$.  Let
$\mathfrak g=\mathfrak{gl}_2$, with
\[
 B(x,y)=\operatorname{tr}(xy)-\tfrac12\operatorname{tr}(x)
                                      \operatorname{tr}(y),
 \qquad P=\sum_{i,j}E_{ij}\otimes E_{ji}.
\]
The identity matrix is retained; it is central and lies in the radical
of $B$.  Fix a basis $E_a$ and write $[E_a,E_b]=f_{ab}^cE_c$,
summing repeated colour indices.  The notation $\sum_Px\otimes y$
means contraction with the trace tensor $P$.

Let $A$ be the exterior algebra on $c_n^a$, $n\in\mathbb Z$, with
cochain degree one and weight $-n-1$.  Give $\hbar$ degree zero and
weight $-1$.  For every integer $L$ let $I_L$ be the exterior ideal
generated by modes $n\geq L$, and set
\[
 G=\varprojlim_{L\to\infty}A/I_L,\qquad
 \widehat I_L=\ker(G\to A/I_L).
\]
Each quotient is an ordinary polynomial exterior algebra.  Thus only
finitely many monomials of a completed element survive any fixed upper
mode bound.  This topology is different from a completion by the total
number of generators.

The continuous even loop derivations and translation are
\[
 L_{x,j}c_n^a=-f_{xb}^a c_{n+j}^b,
 \qquad [L_{x,i},L_{y,j}]=L_{[x,y],i+j},
 \qquad \partial c_n^a=-nc_{n-1}^a.
\]
Define the first-order product on $G\otimes R_2$ by
\begin{equation}\label{cyquant:eq:ghost-product}
 f*g=fg+\hbar\mu_1(f,g),\qquad
 \mu_1(f,g)=-\frac12\sum_{r\geq0}\sum_P
            (L_{x,-r-1}f)(L_{y,r}g).
\end{equation}

\begin{proposition}\label{cyquant:prop:ghost-completion}
This is a continuous associative unital product, free over $R_2$.
Every $\widehat I_L\otimes R_2$ is a closed left ideal.  The smooth
left modules
$F_L=(G\otimes R_2)/(\widehat I_L\otimes R_2)$ are jointly faithful;
their weight pieces are finite-dimensional and their weights are bounded
below.
\end{proposition}

\begin{proof}
Modulo $\widehat I_L$, the right input is a finite exterior polynomial.
For sufficiently large $r$, $L_{y,r}$ kills it in that quotient.
After this finite range is fixed, an enlarged bound on the left input
suffices.  This proves convergence and continuity.  The nonnegative
right shifts preserve $\widehat I_L$, proving the left ideal assertion.
For even derivations $U,V$, the map $(f,g)\mapsto(Uf)(Vg)$ is a
Hochschild two-cocycle by the two Leibniz rules.  Thus $\mu_1$ is a
cocycle and the product is associative modulo $\hbar^2$.  Derivations
kill $1$, which remains its unit; the underlying module is unchanged.

Write $\rho_L$ for the left action and $1_L$ for the quotient of $1$.
Then
$\rho_L(f)1_L=f\bmod\widehat I_L$.  The intersection of these
kernels is zero.  For any fixed polynomial vector only finitely many
right shifts occur, so sufficiently large left generator modes kill
it as well.  There are finitely many nonpositive-weight generators in
each quotient and four generators in each sufficiently large positive
weight.  Exterior monomials give the claimed weight bounds.
\end{proof}

These quotients are modules, not quotient algebras for $*$.  For example,
in the basis $H=E_{11}-E_{22}$, $E=E_{12}$, $F=E_{21}$, $I=E_{11}+E_{22}$,
\[
 c_1^E*c_{-1}^H
 =-\hbar(c_{-1}^Hc_0^E+c_0^Hc_{-1}^E)
                 \pmod{\widehat I_1},
\]
although $c_1^E\in I_1$.  This one-sided feature is essential when
taking transposes.

\subsection{Fields and the first curved identities}

Separate jet states from modes.  Let
$B_0=\Lambda(\xi_j^a:j\geq0)$ with
$T\xi_j^a=(j+1)\xi_{j+1}^a$ and $c^a=\xi_0^a$.
Each $\xi_j^a$ has cochain degree one and weight $j$.
Its classical fields are $Y_0(b,z)c=(e^{zT}b)c$.
The coadjoint operators
$\ell_{x,m}\xi_j^a=-f_{xb}^a\xi_{j-m}^b$, with negative jets zero,
extend by the Leibniz rule.  Set
\begin{equation}\label{cyquant:eq:ghost-exchange}
 s_B(z)=\sum_{m,r\geq0}(-1)^m\binom{m+r}{m}z^{-m-r-1}
                     \sum_P\ell_{x,m}\otimes\ell_{y,r}.
\end{equation}
This sum is finite on each state tensor.  The state fields and exchange
to first order are
\[
 Y_B=Y_0(1-\tfrac\hbar2s_B),\qquad S_B=1+\hbar s_B.
\]
They define a quantum vertex superalgebra with the unchanged vacuum
and translation.  Indeed translation of the kernel
$P/(z+t_1-t_2)$ gives $[T\otimes1,s_B]=-\partial_zs_B$,
and its Leibniz rule on a product gives the two hexagon identities.
Both iterated fields have the same correction
\[
 1-\tfrac\hbar2
       \bigl(s_{12}(z_1-z_2)+s_{13}(z_1)+s_{23}(z_2)\bigr),
\]
so classical associativity gives weak associativity.  Reversing the
first two inputs changes the sign of $s_{12}$ and gives the specified
exchange.

For a state $b$, substitute the full mode series
$c^a(z)=\sum_nc_n^az^{-n-1}$ and $T=\partial_z$ in its classical
exterior expression, obtaining coefficients $b_{b,n}\in G$.
At a fixed coefficient, the sum of indices is fixed; an upper bound
on all indices bounds them below as well.  Hence the sum converges in
$G$.  Its field on $F_L$ is
\[
 Y_L(b,z)=\sum_n\rho_L(b_{b,n})z^{-n-1}.
\]
The covariance identity
\[
 L_{x,j}b_b(z)=\sum_{m\geq0}\binom jm z^{j-m}b_{\ell_{x,m}b}(z)
\]
follows first on jets, then on products.  At $j=-r-1$ it turns
\eqref{cyquant:eq:ghost-product} into precisely the half-twisted field
formula, proving the module identities for every state.  The exchange
is nonidentity: $(S_B-1)(c^E\otimes c^E)=2\hbar z^{-1}c^E\otimes c^E$.

Fix a level $K\in\mathbb C$.  On the classical exterior algebra define
odd derivations
\begin{equation}\label{cyquant:eq:ghost-D}
 \begin{aligned}
 D_0c_n^c&=-\tfrac12f_{ab}^c\sum_{p+q=n-1}c_p^ac_q^b,\\
 D_1&=-\tfrac K2\sum_{r\geq0,b}(r+1)
       \bigl(c_{-r-2}^bL_{E_b,r}+c_r^bL_{E_b,-r-2}\bigr),
 \end{aligned}
\end{equation}
and, writing $B_{ab}=B(E_a,E_b)$, curvature coefficients
\begin{equation}\label{cyquant:eq:ghost-chi}
 \chi_0=-K\sum_{r\geq0,a,b}B_{ab}(r+1)c_{-r-2}^ac_r^b,
 \qquad
 \chi_1=K^2\sum_{r\geq0,a,b}B_{ab}\binom{r+2}{2}c_{-r-3}^bc_r^a.
\end{equation}
For a fixed output bound $L$, the first two maps need only input
bounds $2L$ and $2L+1$ respectively.  The curvature sums converge
because their right factors have index $r$.  Thus all formulas define
continuous maps in the chosen completion.

\begin{theorem}\label{cyquant:thm:ghost-curved-first}
The operator $D=D_0+\hbar D_1$ is an odd derivation of $*$ and, with
$\chi=\chi_0+\hbar\chi_1$, satisfies
\[
 D^2=[\chi,-]_*,\qquad D\chi=0.
\]
For $K\ne0$, one has $D^2\ne0$.
\end{theorem}

\begin{proof}
Pair $c_n^a$ with $E_at^{-n-1}$.  Then $D_0$ is the loop
Chevalley--Eilenberg differential, so $D_0^2=0$ and
$[D_0,L_{x,j}]=0$.  It differentiates each term of $*$;
$\hbar D_1$ is a classical derivation at its order.
For the contraction $\iota_{x,j}c_n^a=x^a\delta_{n,-j-1}$,
Cartan's identity gives $[D_0,\iota_{x,j}]=L_{x,j}$.
The cochain $\chi_0$ evaluates to
$-KjB(x,y)\delta_{j+\ell,0}$ on $(xt^j,yt^\ell)$.
Invariance of $B$ makes it a cocycle.  Thus, with
$x_n^\flat=B(x,E_a)c_n^a$,
\[
 D_0\chi_0=0,\qquad
 \iota_{x,j}\chi_0=-Kj x_{j-1}^\flat,
 \qquad L_{x,j}\chi_0=-KjD_0x_{j-1}^\flat.
\]
Substitution in the antisymmetrization of $\mu_1(\chi_0,-)$ gives
$[D_0,D_1]$.  The contraction of $B$ with $P$ is the traceless
projection; the missing scalar acts by the zero adjoint operator, so
no trace current is removed.  This proves the curvature commutator.
Reindexing the same formula gives $D_1\chi_0=-D_0\chi_1$, proving
the Bianchi identity.  All rearrangements are finite after an output
cutoff.  For nonvanishing, the coefficient of
$c_{-6}^Ec_1^Hc_1^F$ in $D^2c_{-1}^H$ is $-\hbar K$.
\end{proof}

The classical zero mode of $w_K=\frac K2B_{ab}c^aTc^b$ is $\chi_0$.
The additional term $\chi_1$ is a convergent mode operator.  For
$K\ne0$ it is not the zero mode of a finite quadratic local state:
on the fixed mode-sum line its antisymmetrized coefficient is
proportional to
$(q+1)(q+2)\operatorname{sgn}(q+\frac32)/2$, whereas a finite
quadratic differential polynomial has a polynomial coefficient in
$q$.  The two half-lines cannot be values of that same polynomial.
The common completion permits the operator without turning it into
such a local state.

\subsection{The affine central term from an auxiliary generator}

Adjoin an odd generator $\tau$ of weight zero with
\[
 [\tau,g]_{\mathrm s}=D_0g,\qquad \tau^2=\chi_0.
\]
Its normal form is $G\oplus G\tau$: the overlapping reductions agree
by $D_0^2=0$, $D_0\chi_0=0$, and classical centrality of $\chi_0$.
Extend the loop derivations by
\[
 \widetilde L_{x,j}\tau=\alpha_{x,j}:=-Kj x_{j-1}^\flat.
\]
The identities above preserve both relations, and invariance gives
$L_{x,i}\alpha_{y,j}-L_{y,j}\alpha_{x,i}=\alpha_{[x,y],i+j}$.
Thus these are actual loop derivations.  Twisting their product by
the same formula as \eqref{cyquant:eq:ghost-product} gives
\[
 \tau\star c+c\star\tau=(D_0+\hbar D_1)c,
 \qquad \tau\star\tau=\chi_0+\hbar\chi_1.
\]
The auxiliary generator records the curved data; it is not an extra
polarized ghost state.

Let $z$ now denote the central homogenizing generator dual to $\tau$
in this quadratic presentation, rather than an insertion coordinate.
Its dual relations are
\[
 [J_x(m),J_y(n)]
  =zJ_{[x,y]}(m+n)+KmB(x,y)\delta_{m+n,0}z^2,
 \qquad [z,J_x(m)]=0.
\]
The negative transpose of $\widetilde L$ is
\[
 \delta_{x,j}J_y(n)
   =J_{[x,y]}(j+n)+KjB(x,y)\delta_{j+n,0}z,
 \qquad \delta_{x,j}z=0.
\]
The central summand is the pairing with $\widetilde L_{x,j}\tau$.
After $z=1$, this is the affine inner derivation $[J_x(j),-]$.
For $j\geq0$ it descends to the vacuum module and gives
\begin{equation}\label{cyquant:eq:affine-vacuum}
 \delta_{x,1}(J_y(-1)\mathbf1)=KB(x,y)\mathbf1.
\end{equation}
Omitting the auxiliary image would lose this vacuum component.  The
finite quadratic transpose consequently retains the level and scalar
affine relation before taking the vacuum quotient.

\section{The second-order product and its exact range}
\label{cyquant:sec:second-order}

The second coefficient of the product must correct associativity and
remain compatible with translation.  Write
\[
 \mathcal K=\sum_{r\geq0}\sum_P L_{x,-r-1}\otimes L_{y,r},
 \qquad \kappa_2=\tfrac18\mu_0\mathcal K^2,
\]
where $\mu_0$ is the classical exterior product.
The defect left after $\kappa_2$ is a triderivation.  On generators
it is
\[
 -\frac18\sum_{\substack{i+j+k=-2\\i<0,\ k\geq0}}
       \theta^{xyz}(L_{x,i}u)(L_{y,j}v)(L_{z,k}w),
 \qquad \theta=[P_{12},P_{23}].
\]
An explicit correction will cancel it by its three-mode shifts.

For $\xi\in\mathfrak g^*$ write $\xi_n=\xi(E_a)c_n^a$.
Let $\pi$ be the projection onto the annihilator of $I$, and let
$\Omega_{ab}$ be the trace Casimir acting coadjointly in factors
$a,b$.  Equivalently $\pi=\frac14\sum_Pxy$ in the coadjoint
representation.  Put
$\Sigma_{ab}=\pi_a\pi_b+\Omega_{ab}/2$ and
$\{A,B\}=AB+BA$.  Define six colour operators
\begin{equation}\label{cyquant:eq:three-mode-colours}
 \begin{aligned}
 X_{12}&=\tfrac16\Sigma_{23}\Omega_{12},&
 X_{13}&=\tfrac16\Sigma_{23}\Omega_{13},\\
 X_1&=\tfrac16\pi_1\Omega_{23}
       +\tfrac1{24}\bigl(\{\Omega_{12}+\Omega_{13},\Omega_{23}\}
                         -\{\Omega_{12},\Omega_{13}\}\bigr),\\
 Y_1&=\tfrac16\Sigma_{12}\Omega_{23},&
 Y_2&=\tfrac16\Sigma_{12}\Omega_{13},\\
 Y_{12}&=\tfrac16\pi_3\Omega_{12}
       +\tfrac1{24}\bigl(\{\Omega_{13}+\Omega_{23},\Omega_{12}\}
                         -\{\Omega_{13},\Omega_{23}\}\bigr).
 \end{aligned}
\end{equation}
For shifts $e_1+e_2+e_3=-2$, let $J(e)=\{i:e_i<0\}$.
It is a nonempty proper subset of $\{1,2,3\}$.
Set $X(e)=X_{J(e)}$ when $1\in J(e)$ and zero otherwise;
set $Y(e)=Y_{J(e)}$ when $3\notin J(e)$ and zero otherwise.
For a kernel $Z$, define
\[
 \Psi_Z((\xi_1)_{n_1},(\xi_2)_{n_2},(\xi_3)_{n_3})
   =\sum_{e_1+e_2+e_3=-2}
           \bigl(Z(e)(\xi_1\otimes\xi_2\otimes\xi_3)\bigr)_{n+e},
\]
where $(\eta_1\otimes\eta_2\otimes\eta_3)_q
=(\eta_1)_{q_1}(\eta_2)_{q_2}(\eta_3)_{q_3}$.
Only finitely many terms survive each upper bound.

Define $\beta(u;v\wedge w)=\Psi_X(u,v,w)$ and
$\beta(u\wedge v;w)=\Psi_Y(u,v,w)$.  To extend it to ordered
exterior monomials, choose each subset of one or two factors once.
For factorizations $f=\varepsilon uf'$, $g=\varepsilon'vwg'$ add
$\varepsilon\varepsilon'\beta(u;v\wedge w)f'g'$.
For $f=\varepsilon uvf'$, $g=\varepsilon'wg'$ add
$\varepsilon\varepsilon'(-1)^{|f'|}\beta(u\wedge v;w)f'g'$.
Here $\varepsilon,\varepsilon'$ are the signs of moving the selected
factors to the front.  The colour symmetries make each selected pair
alternating, so this defines an exterior operation.

For an equivariant colour operator $A=\sum_iX^i\otimes Y^i$, let
$X^i_j\xi_n=(X^i\xi)_{n+j}$ be its extension as an even derivation,
and put
\[
 \pi_A(f,g)=\sum_{r\geq0,i}(r+1)
                 (X^i_{-r-2}f)(Y^i_rg).
\]
On generators this is the expansion of $A/(z-w)^2$.
If $\tau_{12}$ denotes exchange of the two colour factors, set
\begin{equation}\label{cyquant:eq:level-colour}
 A_*=\tfrac12(1-\tau_{12})(\pi\otimes\pi)-\tfrac14\Omega,
 \qquad A_2=\tfrac12(1+\tau_{12})(\pi\otimes\pi),
 \qquad A_K=A_*+(K+2)A_2.
\end{equation}
The subscript on $A_2$ here names this colour operator, not a precision
quotient.  Define
\begin{equation}\label{cyquant:eq:second-product}
 f*_Kg=fg+\hbar\mu_1(f,g)+\hbar^2\mu_2^{(K)}(f,g),
 \qquad \mu_2^{(K)}=\kappa_2+\beta+\pi_{A_K}.
\end{equation}

\begin{theorem}\label{cyquant:thm:second-product}
Over $R_3$, the product \eqref{cyquant:eq:second-product} is
continuous, associative, and unital.  Translation $\partial$ and the
constant coadjoint operators are derivations.  Each
$\widehat I_L\otimes R_3$ is a closed left ideal, and its smooth
quotient modules are jointly faithful.  The polarized quotient
$B_3=(G\otimes R_3)/(\widehat I_0\otimes R_3)$ has a quantum vertex
superalgebra structure with vacuum $1_0$, translation $\partial$,
and exchange $\exp(\hbar s_B)$.  On every quotient $F_L$, each
composite state acts by its classical exterior coefficients through
the product $*_K$.
\end{theorem}

\begin{proof}
In the colour representation,
\[
 [\Omega_{12},\Omega_{23}]
 =[\Omega_{23},\Omega_{13}]=[\Omega_{13},\Omega_{12}].
\]
The operator $\Sigma_{ab}$ has image symmetric in factors $a,b$:
on their antisymmetric traceless summand the Casimir is $-2$, so
$\pi_a\pi_b+\Omega_{ab}/2$ vanishes.  A shift kernel constant on
one half-plane, with image symmetric in the two remaining factors,
acts by zero after exterior multiplication.  Indeed the finite set
of surviving mode pairs is stable under exchange and its terms cancel
in pairs, including zero at a fixed pair.

The colour identities give
\[
 X_1-Y_1+\tfrac18\theta
 =X_{12}-Y_{12}+\tfrac18\theta
 =X_{13}-Y_2=\tfrac1{12}(\Omega_{23}-\Omega_{12})\Omega_{13}.
\]
On the six sign regions these turn the remaining associator into two
of the vanishing half-plane kernels just described.  The extension
rule for $\beta$ makes its Hochschild coboundary a triderivation:
terms selecting factors from only two inputs cancel in pairs, leaving
one selected factor from each input.  The generator computation
therefore proves associativity on all monomials.  The additional
$\pi_{A_K}$ is a biderivation and hence a Hochschild cocycle.

The same half-plane cancellation proves
$\beta(G,\widehat I_L)\subset\widehat I_L$.
For a right-input polynomial with least mode $p_0<L$, the bound
$M=3L-2p_0+2$ gives
$\beta(\widehat I_M,g)\subset\widehat I_L$.
For the one-plus-two part this follows from the sum of its three
output modes; for the two-plus-one part the possible low output lies
in a symmetric half-plane and cancels.  The other product terms
satisfy the same bound because their right shifts are nonnegative.
These estimates prove continuity, the left ideals, smoothness, and
the extension of associativity to the completion.  Evaluation at $1_L$
again separates coefficients.

For translation, a kernel depending only on $J(e)$ has defect
$\sum_i(e_i'+1)Z(e'+\delta_i)$ on $\sum e_i'=-3$.
All terms with nonzero coefficient use the same sign region, and
their coefficients sum to zero.  This proves translation for $\beta$.
For the double-pole kernel the boundary term is
$(a-1)(a-2)[a=2]A_K=0$.  Thus $\pi_{A_K}$ also commutes with
translation.  Equivariance is built into the colour operators.

Finally $\tau_{12}A_K\tau_{12}=A_K$.  The two expansions of its
double pole therefore agree after clearing the diagonal, giving the
generator exchange.  Reconstruction from the creating generating
fields gives the state algebra.  For a fixed module vector with least
mode $p_0$, use the same bound $M=3L-2p_0+2$ before taking an
iterate.  A word of $r$ fields with total derivative order $d$ has
no power below $-rM-d$, so one power of the associativity denominator
suffices independently of output weight.  Taylor expansion identifies
the iterate with the field of the composite state.  This proves the
all-state module assertion and extends exchange from generators to
all states.
\end{proof}

\subsection{The derivation and curvature at precision three}

Let $\delta:\mathfrak g^*\to\mathfrak g^*\otimes\mathfrak g^*$ be
$\delta c^c=-\frac12f_{ab}^cc^a\otimes c^b$.
Put
\[
 \Delta_K=[D_0,\mu_2^{(K)}]+[D_1,\mu_1].
\]
These brackets denote the failure to be a derivation, with the odd
Leibniz sign.  Write $\delta(z-w)=\sum_{n\in\mathbb Z}z^{-n-1}w^n$
for the formal delta function, and let a superscript $+$ mean the sum
of a bilinear expression and its input reversal on odd generators.
The diagonal symmetrization vanishes:
\begin{equation}\label{cyquant:eq:diagonal-cancellation}
 \Delta_K^+
  =-(1+K/2)\delta(z-w)S_1+\tfrac{K+2}{2}\delta(z-w)S_1=0,
\end{equation}
where
$S_1(\xi,\eta)=(\partial_1(\pi\xi\otimes\delta\eta+
\pi\eta\otimes\delta\xi))(w,w,w)$.
The first term comes from the internal collisions in $\beta$ and
the level term $D_1$; the second is $[D_0,\pi_{(K+2)A_2}]^+$.
The identity follows by replacing the difference of the two
double-pole expansions by $\partial_w\delta(z-w)$ and taking the
first diagonal Taylor coefficient.

Order generators by mode and then colour.  There is a unique
continuous odd map vanishing on generators that corrects this defect:
\begin{equation}\label{cyquant:eq:Dtwo}
 D_2(u_1\cdots u_k)=
 -\sum_{i=1}^{k-1}(-1)^{i-1}u_1\cdots u_{i-1}
                          \Delta_K(u_i,u_{i+1}\cdots u_k).
\end{equation}
To verify well-definedness, exchange two neighbouring odd generators.
The odd Hochschild cocycle identity for $\Delta_K$ reduces the
difference to their symmetrization, which is zero by
\eqref{cyquant:eq:diagonal-cancellation}.  The same argument handles
a repeated generator.  The first term of the recursion then proves
the derivation identity on a generator and a monomial, and the
cocycle identity extends it to every pair.  The upper mode estimates
extend it continuously.  The difference of two such maps would be a
continuous classical derivation zero on all generators, proving
uniqueness.

Define the classical derivation $E_2$ and the next curvature by
\begin{equation}\label{cyquant:eq:second-curvature}
 \begin{aligned}
 E_2\xi_n&=-\frac{3K+2}{2}
       \sum_{p+q=n-3}(p+1)(p+2)(\delta\xi)_{p,q},\\
 \chi_2&=\frac{K(K^2+9K+6)}2
       \sum_{r\geq0,a,b}\binom{r+3}{3}B_{ab}c_{-r-4}^ac_r^b,\\
 D&=D_0+\hbar D_1+\hbar^2(D_2+E_2),\qquad
 \chi=\chi_0+\hbar\chi_1+\hbar^2\chi_2.
 \end{aligned}
\end{equation}
The map $D$ is a derivation of $*_K$.  The curved identities will
follow from a current pairing in every weight, rather than from a
finite weight window.

Let $V=V^{KB}(\mathfrak{gl}_2)$, with currents $J_x$ satisfying
$Y(J_x,z)J_y\sim KB(x,y)z^{-2}\mathbf1+z^{-1}J_{[x,y]}$.
For trace-dual bases $e_a,e^a$ of $\mathfrak{sl}_2$, put
\[
 \begin{aligned}
 s_A(z)&=\sum_{m,n\geq0}(-1)^m\binom{m+n}{m}z^{-m-n-1}
                          \sum_a e_a(m)\otimes e^a(n),\\
 M_b^p&=\sum_{a,n\geq0}e_a(-n-1)[e^a,e_b](n+p),\\
 a_0(z)&=\sum_{p,q\geq0}(-1)^p\binom{p+q}{p}z^{-p-q-1}
                          \sum_bM_b^p\otimes e^b(q),\qquad
 \check a_0(z)=\tau_{12}a_0(-z)\tau_{12}.
 \end{aligned}
\]
Each sum is finite on a fixed pair of states.  The polar fields at
this precision are
\begin{equation}\label{cyquant:eq:polar-current}
 Y^P=Y\left(1+\tfrac\hbar2s_A+
       \hbar^2\bigl(s_A^2/8+(a_0+\check a_0)/24\bigr)\right),
 \qquad S_A=\exp(-\hbar s_A).
\end{equation}
The diagonal completion $a_0$ compensates the second-order
associator of the half twist.  The loop commutator, its primitive
coproduct, and the rational Casimir identity give the corresponding
associativity and exchange identities coefficientwise.  The scalar
current is an unchanged commutative differential-polynomial factor:
all its nonnegative modes vanish on $V$.

Define equivariant weight-three states for $x\in\mathfrak{sl}_2$ by
\[
 \Psi_1(x)=\sum_a[e_a,x](-2)e^a,
 \qquad
 \Psi_2(x)=\sum_a\operatorname{sym}(e_a(-1)e^a(-1)x),
\]
where $\operatorname{sym}$ averages the six orders.  Set
\begin{equation}\label{cyquant:eq:current-change}
 \theta_K(x)=-\frac{3K^2+33K+22}{72}T^2x
                  +\frac38\Psi_1(x)-\frac16\Psi_2(x),
 \quad \theta_K(I)=0,
 \quad \widehat J_x=J_x+\hbar^2\theta_K(x).
\end{equation}
This is an invertible change on the same polar state space.

\begin{theorem}\label{cyquant:thm:all-mode-mc}
For every level $K$, smooth polar current module $M$, and polarized
ghost module $F_L$ with $L\geq0$, put
\[
 Q_{M,L}=\sum_{a,j}\widehat J_a(j)\otimes\rho_L(c^a_{-j-1}).
\]
Then, modulo $\hbar^3$,
\begin{equation}\label{cyquant:eq:all-mode-mc}
 Q_{M,L}^2+\mathcal DQ_{M,L}+1\otimes\rho_L(\chi)=0,
 \qquad D^2=[\chi,-]_{*_K},\qquad D\chi=0.
\end{equation}
Here $\mathcal D$ differentiates the completed ghost coefficients
before their action on $F_L$; no descent of $D$ to an arbitrary
quotient is assumed.
\end{theorem}

\begin{proof}
We give the residue and exchange terms in the tensor state space
$W=V^P\otimes B_3$.  Let $T$ be its total translation and put
$d_{\mathrm{CE}}c^c=-\frac12f_{ab}^cc^ac^b$.
Define $\tau=\sum_aJ_a\otimes c^a$,
$\Gamma=(1\otimes d_{\mathrm{CE}})\tau$,
$\widehat\tau=\sum_a\widehat J_a\otimes c^a$, and let $T_A$ act on
the current factor.  Define
\[
 \begin{aligned}
 e(c^c)&=\tfrac{3K+2}{4}f_{ab}^c(T^2c^a)c^b,
 &\zeta&=\sum_cJ_c\otimes e(c^c),\\
 v_K&=-\tfrac{K(K^2+9K+6)}{24}B_{ab}c^aT^3c^b,
 &w_K&=\tfrac K2B_{ab}c^aTc^b.
 \end{aligned}
\]
The classical modes of $e$ are $E_2$ and the zero modes of $w_K,v_K$
are $\chi_0,\chi_2$.  Direct residue expansion of
\eqref{cyquant:eq:polar-current} and the ghost generator field gives
\begin{equation}\label{cyquant:eq:mc-primitive}
 \tfrac12\widehat\tau_{(0)}\widehat\tau
 +(1\otimes d_{\mathrm{CE}})\widehat\tau+w_K
 =\tfrac{\hbar K}{2}T(\Gamma+w_K)
  +\hbar^2\left(-\zeta-v_K+\tfrac{3K+2}{4}TT_A\Gamma\right).
\end{equation}
This is a finite state calculation.  More explicitly, write
$[\hbar^i]Y^P(J_a,z)J_b=\sum_\ell z^\ell P_{i,\ell}^{ab}$ and
$C_0=1$, $C_1=-\Omega/2$, $C_2=\Omega^2/8+A_K$.
The order-$r$ residue of $\tau_{(0)}\tau$ is
\[
 \sum_{i+j=r}\sum_{a,b,u,v}\sum_{\ell=-i-2}^{j-1}
       (C_j)_{ab}^{uv}P_{i,\ell}^{ab}
                        \otimes\xi_{j-\ell-1}^u\xi_0^v.
\]
The sixteen current pairs and these finite intervals, together with
the two linear insertions of $\theta_K$, give respectively the
coefficients $0$, $\frac K2T(\Gamma+w_K)$, and the last coefficient
in \eqref{cyquant:eq:mc-primitive}.  The cubic and quadratic current
terms cancel.  Thus the local residue in the desired operator identity
is a total translation and has zero zero mode.

There are also exchange contributions.  Put
$R(x)=\sum_aJ_{[x,E_a]}\otimes c^a$,
$\vartheta(x)=1\otimes x^\flat$, and
\[
 \mathcal R=\sum_P(R(x)\otimes\vartheta(y)
                   -\vartheta(x)\otimes R(y)),\qquad
 \mathcal H=\sum_P\vartheta(x)\otimes\vartheta(y).
\]
The tensor exchange is $S_W=\exp(-\hbar s_W)$ for $s_W=s_A-s_B$.
Invariance of the canonical tensor gives
\[
 s_W(z)\tau^{\otimes2}=Kz^{-2}\mathcal R-2K^2z^{-3}\mathcal H,
 \qquad
 s_W(z)^2\tau^{\otimes2}=2Kz^{-3}\mathcal R-8K^2z^{-4}\mathcal H.
\]
The second equation uses the ghost Casimir eigenvalues $-2,-4$ on
the adjoint and invariant colour summands.
For a rational tensor term $z^{-p}u\otimes v$, the correction to a
mode commutator is computed by
\[
 \mathfrak B_{m,n}(z^{-p};u,v)
   =(-1)^p\sum_{r\geq0}\binom{p+r-1}{r}
                          v_{(n-p-r)}u_{(m+r)},
\]
with an additional $(-1)^{|u||v|}$ for the tensor term.
The rightmost mode makes the sum finite on each smooth vector.
Expanding $S_W^{-1}$ gives both
$\hbar\mathfrak B(s_W)$ and $\hbar^2\mathfrak B(s_W^2)/2$.
In the first term the reversed product must itself be retained
through first order.

Its classical part cancels $D_1$ and $\chi_1$.
Its first-order part has coefficient
\[
 -\hbar^2\left(\frac K2\sum_{c,j}J_c(j)\otimes U^c_{-j-1}
                       +2K^2\,1\otimes N_3\right),
\]
where
\[
 \begin{aligned}
 U_n^c&=\sum_{a,b,r\geq0}\binom{r+2}{2}f_{ba}^c
       (c_{-r-3}^bc_{n+r}^a-c_{n-r-3}^ac_r^b),\\
 N_3&=\sum_{a,b,r\geq0}\binom{r+3}{3}B_{ab}c_{-r-4}^bc_r^a.
 \end{aligned}
\]
The second exchange term is its negative.  The coefficients follow
from the same two Casimir eigenvalues and
$\sum_{r+t=u}(r+1)=\binom{u+2}{2}$,
$\sum_{r+t=u}\binom{r+2}{2}=\binom{u+3}{3}$.
This proves the first identity in \eqref{cyquant:eq:all-mode-mc}.
Every rearrangement is finite after the ghost upper bound and the
smooth current bound on a fixed vector.

For the curved identities let
$\mathcal E=D^2-[\chi,-]=\hbar^2\mathcal E_2$ and
$D\chi=\hbar^2\mathcal C_2$.  The first is a continuous classical
derivation.  Apply the operator identity to the affine vacuum module
and all $1_L$; their separation gives an identity of completed
coefficients.  Its Bianchi identity is
\[
 \sum_{a,j}\widehat J_a(j)\otimes\mathcal E(c^a_{-j-1})
                        +1\otimes D\chi=0.
\]
The trace ghost is killed by $D$ and is supercentral for $*_K$:
every pair or triple correction has a coadjoint action or traceless
projection on that input.  Fixing a cubic ghost monomial of mode sum
$s$, the remaining coefficient is $J_x(j)+\lambda1=0$ on the
classical affine vacuum module, with $x\in\mathfrak{sl}_2$ and
$j=-s-5$.  For $j<0$ the vacuum and PBW give $x=0$.
For $j\geq0$ the vacuum first gives $\lambda=0$; applying to
$J_y(-j-1)\mathbf1$ gives $J_{[x,y]}(-1)\mathbf1=0$.
The central term is absent because the indices sum to $-1$.
PBW and the zero centre of $\mathfrak{sl}_2$ imply $x=0$ at every
level.  Thus $\mathcal E_2$ vanishes on all generators and
$\mathcal C_2=0$.  Continuity proves the two full identities.
\end{proof}

Transposing the quadratic product, differential, and curvature on
finite perfect coefficient blocks now gives a current algebra at
each precision $1\leq N\leq3$.  Homogenize the linear and scalar
terms with the generator dual to $\tau$, set that generator to one,
and take the quotient killing nonnegative current modes on the vacuum.
The map to the corrected polar currents satisfies every relation by
\eqref{cyquant:eq:all-mode-mc}.  It is an isomorphism on the singleton
vacuum module.  Indeed the classical affine relations reorder every
word, with an error divisible by $\hbar$; iterating this process
terminates because $\hbar^N=0$.  Thus ordered PBW words span the
source.  Their images are the PBW basis of the flat polar module,
which proves independence without assuming source flatness.
The transposed singleton isomorphism preserves translation, the
outgoing vacuum, and every ordered-word coefficient, including the
trace direction.  Native binary germs and their supported residue
maps require the additional comparison described below.

\subsection{The integral current comparison retains the trace}

The polar current algebra also has a field comparison with the integral
quantum affine algebra; its trace correction differs from the change
\eqref{cyquant:eq:current-change} used in the ghost pairing.
Let $\mathcal U_K$ be the full matrix quantum affine vacuum algebra
with creation matrix $L(v)=I-\hbar j(v)$ and normalized rational
matrix
\[
 R(w)=f(\hbar/w)(I-\hbar P/w),\qquad
 f(t)=1+\tfrac12t+\tfrac58t^2+\tfrac{11}{16}t^3
                         +\tfrac{83}{128}t^4+O(t^5).
\]
The scalar normalization and the integral creation lattice are those
of the rational construction in Section~\ref{cyrep:sec:rtt-double};
here the full trace-current sector is retained.  Write $j=j^{[0]}$,
$Tj^{[n]}=(n+1)j^{[n+1]}$, and interpret products of entries of
$j^{[n]}$ as products in the creation algebra acting on the vacuum.
They are not being identified with quantum $(-1)$-products.

Put $u=J_I$ and $C=\sum_a e_a(-1)e^a$.
For the polar completion $a_\mu=a_0+\mu\partial_zs_A$ define
\[
 \begin{aligned}
 \Theta_\mu(x)={}&
 \left(-\frac{K^2}{12}-\frac K{24}+\frac{13}{72}
                         +\frac{\mu K}{24}\right)T^2x
 +\left(\frac{K+1}{8}-\frac\mu{24}\right)\Psi_1(x)
 +\frac16\Psi_2(x)\\
 &+\frac12u(-1)Tx+\frac14u(-1)u(-1)x,
       \qquad x\in\mathfrak{sl}_2,\\
 \Theta_\mu(I)={}&\frac12u(-1)C.
 \end{aligned}
\]
The map on currents
\begin{equation}\label{cyquant:eq:integral-current-map}
 J_{E_{ij}}\longmapsto
 j_{ij}+\frac\hbar2\sum_a j_{ia}j_{aj}
                   +\hbar^2\widetilde{\Theta_\mu(E_{ij})}
\end{equation}
extends to a graded, equivariant isomorphism of quantum vertex
algebras $V^P_\mu\simeq\mathcal U_K/\hbar^3\mathcal U_K$.
The tilde is any lift of the classical PBW state; multiplication by
$\hbar^2$ makes it independent of that choice.  This comparison
fixes the vacuum, commutes with translation, and preserves all fields
and exchanges without localizing the level.

Here is how the second coefficient is detected.  For a current change
$\theta$, its singular-product coboundary is
\[
 q_\theta(x,y;z)=\operatorname{Sing}_z
       \bigl(Y(\theta(x),z)y+Y(x,z)\theta(y)\bigr)
                         -z^{-1}\theta([x,y]).
\]
The residual of the first logarithmic change
$j\mapsto j+\hbar j^2/2$ on all sixteen matrix-current pairs is
$-q_{\Theta_\mu}$.  It follows by expanding the mixed RTT relation
through the required scalar order, straightening the creation words,
and using the affine commutator for the polar coefficients.
Both recursions terminate: the classical commutator shortens a word,
and its quadratic correction raises the power of $\hbar$.
The exchange residual after the first coordinate change is
$s_A^2/2$.  Thus \eqref{cyquant:eq:integral-current-map} matches the
complete generating products and exchange.  Translation-compatible
PBW extension and weak associativity then give the equality on every
state.  Flatness and the identity classical reduction give the inverse
by a finite $\hbar$-adic recursion.

In particular, for $y^\circ=y-\operatorname{tr}(y)I/2$, the uncorrected
trace residual is
\[
 q(I,y;z)=-(K+2)
       \left(\frac{u(-1)y^\circ}{z^2}
                      +\frac{T(u(-1)y^\circ)}z\right),
 \qquad q(y,I;z)=-(K+2)\frac{u(-1)y^\circ}{z^2}.
\]
The Casimir identity
$\operatorname{Sing}_zY(C,z)y^\circ
 =2(K+2)(z^{-2}y^\circ+z^{-1}Ty^\circ)$
shows that $I\mapsto u(-1)C/2$ removes exactly these terms.
It is essential away from $K=-2$.  A change only on traceless
currents would leave the trace field comparison false.

\subsection{The next associative order}

The specified precision-three product has a genuine obstruction at
noncritical level.  We state its admissible class before the conclusion.
Let $\nu$ be a continuous even biderivation of the classical exterior
product, of degree zero and weight two, equivariant for the constant
coadjoint action.  Require its symmetrization on two generating fields
to be supported on their diagonal.  Allow a continuous odd first-order
correction $D'$ of degree one and weight one with
$[D_0,D']=[D_0,D_1]$.

\begin{proposition}\label{cyquant:prop:third-obstruction}
Suppose $D_0+\hbar D'+\hbar^2D_2'$ is a derivation of
$\mu_0+\hbar\mu_1+\hbar^2(\kappa_2+\beta+\pi_{A_*}+\nu)$
modulo $\hbar^3$.  If this product has any parity-preserving bilinear
associative correction at order $\hbar^3$, then $K=-2$.
In particular the fixed product $*_K$ has no associative extension
modulo $\hbar^4$ when $K\ne-2$.
\end{proposition}

\begin{proof}
Pass to the quotient by the closed ideal of trace ghosts.  Equivariance
and diagonal support express the relevant symmetric colour coefficient
as $Y=y_1P_1+y_3P_3+y_5P_5$, where $P_i$ project onto the
dimensions $1,3,5$ of the tensor square of the adjoint representation.
The correction $F=D'-D_1$ is a $D_0$-closed odd derivation.
Its finitely many coefficients used here may be averaged over $SU(2)$,
so it may be taken equivariant.  Its remaining component has the form
$F(\xi_n)=\sum_{p+q=n-2}a(p,q)(\delta\xi)_{p,q}$ modulo trace
ghosts, with $a(p,q)=a(q,p)$.

Evaluate the second-order derivation defect on $u=c_{-1}^H$ twice
and the third-order associator on $u$ three times.  Their unknown
correction terms vanish: $u^2=0$, while $D_2'u$ and $\mu_3(u,u)$
are even and commute classically with $u$.  Write $h,e,f$ for the
three traceless ghost colours and take the coefficients of
$m_2=e_{-2}f_{-2}h_{-1}$, $m_3=h_{-3}e_{-1}f_{-1}$, and
$m_4=f_{-4}h_{-1}e_{-1}$.  The $[D_0,F](h_{-2})$ coefficient
at $h_{-2}e_{-2}f_{-1}$ is
$2(a(-3,-1)-a(-2,-1))$.
The $m_2,m_3$ coefficients of $[D_0,\nu](u,u)$ are
$(-y_1+3y_3+y_5)/3$ and $-(y_1+2y_5)/3$; the corresponding
$[F,\mu_1]$ coefficients are zero and
$2(a(-2,-1)-a(-3,-1))$.
The third-order $m_4$ coefficient is
$(2y_1-3y_3+y_5)/12$.  Together with the base level contribution
$K+2$ at $m_3$, these yield
\[
 \begin{aligned}
 a(-3,-1)&=a(-2,-1),\\
 -y_1+3y_3+y_5&=0,\\
 K+2-(y_1+2y_5)/3&=0,\\
 2y_1-3y_3+y_5&=0.
 \end{aligned}
\]
They can be computed from the double-pole expansion and the three
Casimir eigenvalues $-4,-2,2$; the first equality removes the possible
contribution of $F$ to the third line.  Adding the second and fourth
lines gives $y_1+2y_5=0$, and the third gives $K=-2$.
For $*_K$, take $\nu=(K+2)\pi_{A_2}$ and $D'=D_1$.
An arbitrary bilinear extension would have a parity-preserving component
solving the same linear extension equation, so the fixed-product
conclusion excludes it as well.
\end{proof}

At $K=-2$ there is a positive associative extension.  If $a_3$ is
the third-order associator of $*_{-2}$, its total symmetrization on
odd generators is zero.  Substitution of the colour operators into
the rational three-point terms proves their cancellation; the
differences of Laurent expansions are derivatives of diagonal delta
functions, and their contact terms cancel by the same odd exchange
and Casimir identities.  This is an all-mode rational identity.

Here is an explicit continuous primitive of that cocycle.  Order
generators by mode and colour.  For $v>x$ put
$\mathcal S_{vx}(F)=a_3(F,v,x)+a_3(F,x,v)$ and put
$\mathcal S_{xx}(F)=a_3(F,x,x)$.  For $g_1<\cdots<g_k$ set
\[
 \rho_x(g_1\cdots g_k)=
 -\sum_{j:g_j\geq x}(-1)^{k-j}
       \mathcal S_{g_jx}(g_1\cdots g_{j-1})g_{j+1}\cdots g_k,
\]
and define
\begin{equation}\label{cyquant:eq:cubic-primitive}
 \begin{aligned}
 \mu_3(f,g_1\cdots g_k)=\sum_{j=1}^k
 \bigl(&a_3(f,g_1\cdots g_{j-1},g_j)
       -f\rho_{g_j}(g_1\cdots g_{j-1})\\
       &+\rho_{g_j}(fg_1\cdots g_{j-1})\bigr)g_{j+1}\cdots g_k.
 \end{aligned}
\end{equation}
The Hochschild cocycle identity and $\operatorname{Sym}a_3=0$
show that exchanging adjacent generators leaves this formula
consistent with the exterior signs.  Induction on the second word,
then the third input of its coboundary, gives $\delta_0\mu_3=a_3$.
The right ideal and enlarged first-input bounds for $a_3$ give the
same bounds for these recursions, so $\mu_3$ extends continuously and
preserves every right-input upper mode ideal.  Therefore
$*_{-2}+\hbar^3\mu_3$ is associative modulo $\hbar^4$.
The ordering used in this primitive is not preserved by translation.
A translation-compatible field extension and the next curved
differential remain further requirements.  The obstruction above
also retains its stated admissible class; it does not exclude every
possible presentation of a dual field theory.

\section{Finite adjunction before stabilization}
\label{cyquant:sec:native-adjunction}

The singleton comparison does not determine a binary field operation.
That operation uses a relative germ, its quadratic transpose, and the
adjoint in one specified input.  We first construct these maps
classically and then retain the loop cutoff during the first-order
transpose.

Let $V=V^{KB}(\mathfrak{gl}_2)$ and
$V_{\leq a}=\bigoplus_{0\leq d\leq a}V_d$.
The classical ghost germ sends
\[
 \xi_j^\alpha\longmapsto\frac1{j!}\partial_u^j
                         \sum_nc_n^\alpha u^{-n-1}.
\]
Its finite first-input transpose is
\begin{equation}\label{cyquant:eq:native-inserted-action}
 A_x^a(m;u)=\sum_{j=0}^a\binom mj u^{m-j}J_x(j)|_{V_{\leq a}}.
\end{equation}
Negative $m$ occur too: their coefficients pair with the full
annihilation germ before transposition.  On
$\mathscr H_a=\operatorname{Hom}(V_{\leq a},V((u)))$, finite
tensor--Hom adjunction gives
\begin{equation}\label{cyquant:eq:native-hom-action}
 R_x^a(m)F=J_x(m)F-FA_x^a(m;u).
\end{equation}
The minus sign is the antipode action on the dual input.
Let $\epsilon_a(v)=e^{uT}v$ be the creation kernel.
Write $U_K$ for the enveloping algebra of affine mode words with
central generator acting by $K$.

\begin{proposition}\label{cyquant:prop:classical-adjunction}
The cyclic module $\mathscr K_a=U_K\epsilon_a$ is canonically
isomorphic to $V$, by a map $\Psi_a$ with
$\Psi_a(\mathbf1)=\epsilon_a$.  Its inverse is evaluation on the
first vacuum.  These maps commute with restriction of the first-input
bound, and they retain the affine linear and scalar relations.
\end{proposition}

\begin{proof}
Vandermonde's identity gives
$[A_x^a(m),A_y^a(n)]=A_{[x,y]}^a(m+n)$.
The indices in these first-input operators are nonnegative, so their
central cocycle is zero.  Negative precomposition is a representation,
because precomposition reverses order.  Hence
\eqref{cyquant:eq:native-hom-action} satisfies the affine relation,
with its central term occurring once in the output.
For $m\geq0$,
\[
 J_x(m)e^{uT}v=e^{uT}\sum_{j=0}^m\binom mj u^{m-j}J_x(j)v,
\]
so $R_x^a(m)\epsilon_a=0$.  The vacuum quotient supplies $\Psi_a$.
Every $A_x^a(m)$ kills the first vacuum, so evaluation there intertwines
$R_x^a(m)$ with $J_x(m)$ and returns $\mathbf1$ from $\epsilon_a$.
Its composite with $\Psi_a$ is the identity, proving injectivity;
surjectivity holds onto the cyclic module by definition.
Restriction preserves the first-input matrices and the creation kernel.
Smoothness follows from the resulting identification with $V$.
\end{proof}

For output projection $\pi_c:V\to V_{\leq c}$, the resulting
coefficient block is
\[
 \Gamma_{a,b,c}:V_{\leq b}\longrightarrow
 \operatorname{Hom}\bigl(V_{\leq a},
       V_{\leq c}\otimes\langle u^{-a-b},\ldots,u^c\rangle\bigr),
 \qquad \Gamma_{a,b,c}(w)(v)=\pi_c\Psi_a(w)(v).
\]
The range is finite: if the inputs have weights $i,j$ and the output
has weight $k$, the Laurent power is $k-i-j$.
The recursion in negative creation words uses finitely many
precomposition indices, whose total is at most $i$, and finitely many
terms of $e^{uT}$ after output projection.
For at most $d$ additional translations or derivatives, put
$s_0=a+b+c+d+2$.  All required ghost generators lie in
$[-4s_0-1,4s_0-1]$; the quadratic rows can be read with upper bound
$4s_0+1$ and weight interval $[-8s_0-2,8s_0+2]$.
These bounds follow by straightening the finite current words and
then transposing their linear terms.  They provide a common finite
source for each matrix, rather than assuming that a fixed cutoff is
preserved by the differential.

The compatible operation $\mu(v,u)w=\Psi(w)(v)$ satisfies
\[
 \begin{aligned}
 \mu(J_x(-1)\mathbf1,u)J_y(-1)\mathbf1
 ={}&J_y(-1)e^{uT}J_x(-1)\mathbf1\\
 &+u^{-1}e^{uT}J_{[x,y]}(-1)\mathbf1
       +Ku^{-2}B(x,y)\mathbf1.
 \end{aligned}
\]
The same recursion on all creation words identifies it with the
classical affine field, including
$\mu(\mathbf1,u)w=w$, $\mu(v,u)\mathbf1=e^{uT}v$, and its reversed
operation.  Consequently the two ordered coproduct transposes give
separate commuting product comparisons.  The actual residue identity
is equality of
$\operatorname{pp}_u(f\mu(v,u)w)$ with the corresponding field
residue.  For pole order $p$ in $f$, only output weights below
$a+b+p$ contribute.  Spectators and any fixed finite number of normal
derivatives are therefore covered by a common finite enlargement.
This also gives the exchanged $13\mid2$ residue and compatibility
with finite perfect biduality at the classical precision.

\subsection{The two finite terms that cancel}

Let $H$ be the enveloping Hopf algebra of the centrally extended loop
algebra, with primitive generators $x[m]$ and central generator
$\mathbf c$.  Its bracket is
$[x[m],y[n]]=[x,y][m+n]+mB(x,y)\delta_{m+n,0}\mathbf c$.
On ghosts $\mathbf c$ acts by zero, and on currents by $K$.
Write $S$ for the antipode of $H$.
The Hopf counit here belongs to $H$; no character of the fixed-level
affine algebra is required.
For a loop cutoff $M$ put
\[
 r_M=\sum_{0\leq r<M}\sum_Px[-r-1]\otimes y[r],
 \qquad \sigma_M^\pm=1\pm\tfrac\hbar2r_M.
\]
These are counital twists modulo $\hbar^2$, since their linear twist
equations are the sums of the three placements of primitive tensors.
Use $\Delta_\sigma=\sigma\Delta\sigma^{-1}$ and
$m_\sigma=m\sigma^{-1}$.

Twist the auxiliary relative ghost algebra, including its shared
$\tau$, by $\sigma_M^+$.  Its product is the finite sum in
\eqref{cyquant:eq:ghost-product}; commutator with $\tau$ and its
square give the finite $D$ and $\chi$.  The generator-linear relative
germ remains an equivariant quadratic map.  Finite inverse transposition
changes the twist to $\sigma_M^-$, and after setting the generator
dual to $\tau$ to one its relative current map is
\[
 \mathsf J_x(m)\longmapsto
          \mathbb A_x^a(m;u)\otimes1+1\otimes J_x(m)
\]
into the twist of
$(U(\mathfrak g[t]/t^{a+1})\otimes U_K)[u,u^{-1}]$.
Here $\mathbb A$ is \eqref{cyquant:eq:native-inserted-action} with
formal truncated-loop generators in place of their operators.
The whole relative algebra is twisted, including its mixed rows.
The first factor has central level zero and the second has level $K$.

Take the right dual of the finite first-input module
$U=V_{\leq a}\otimes R_2[u,u^{-1}]$, with evaluation in the order
$U\otimes U^{\mathrm r}\to R_2[u,u^{-1}]$.
Its action uses the inverse antipode.  For $\sigma=\sigma_M^-$,
\[
 \mathcal Q_M=m_H(1\otimes S)(\sigma)
      =1+\tfrac\hbar2\sum_{r<M,P}x[-r-1]y[r],
 \qquad S_\sigma^{-1}(h)
   =S^{-1}(\mathcal Q_M)S^{-1}(h)S^{-1}(\mathcal Q_M)^{-1}.
\]
The right-dual comparison contributes precomposition in the order
$A_y^a(r)A_x^a(-r-1)$.  The tensor comparison contributes negative
precomposition on the first factor and the output current on the
second.  Their sum is the actual adjoint comparison
\begin{equation}\label{cyquant:eq:finite-right-dual}
 \begin{aligned}
 \Theta_M(F)={}&F+\tfrac\hbar2\sum_{r<M,P}
                    FA_y^a(r)A_x^a(-r-1)\\
 &-\tfrac\hbar2\sum_{r<M,P}J_y(r)FA_x^a(-r-1)
       =F-\hbar\mathcal C_M^aF,\\
 \mathcal C_M^aF={}&\tfrac12\sum_{r<M,P}(R_y^a(r)F)A_x^a(-r-1).
 \end{aligned}
\end{equation}
The first sum is fixed by right-dual evaluation, not a choice of
normalization.  Its inverse is $1+\hbar\mathcal C_M^a$.

\begin{theorem}\label{cyquant:thm:finite-adjunction}
At every level $K$, the first-order native quadratic germ, with its
full differential and curvature, has a finite adjunction on each
coefficient block.  Its stabilized cyclic kernel is
\begin{equation}\label{cyquant:eq:quantum-kernel}
 \begin{aligned}
 W_u(w)(v)=\Psi(w)(v)+\frac\hbar2
 \sum_{m,n\geq0,P}(-1)^m\binom{m+n}{m}u^{-m-n-1}
          \Psi(J_y(n)w)(J_x(m)v).
 \end{aligned}
\end{equation}
It preserves both cyclic vacua and all native quadratic, linear, and
scalar relations.  Its finite coefficient blocks are compatible with
restriction, inclusion, projection, and finite translation and derivative
enlargements.
\end{theorem}

\begin{proof}
At cutoff $M$, twisting the classical equivariant adjunction leaves
its underlying map $\Psi_a$ unchanged.  Applying the actual right
adjoint comparison \eqref{cyquant:eq:finite-right-dual} gives
$W_{a,M}=(1+\hbar\mathcal C_M^a)\Psi_a$.
The finite twisting functor preserves the module law, so this map
satisfies the full transposed relative relations, including those
containing the differential and curvature.

On $w\in V_{\leq b}$,
$R_y^a(r)\Psi_a(w)=\Psi_a(J_y(r)w)$ vanishes for $r>b$.
Thus $W_{a,M}(w)$ stabilizes as soon as $M>b$, before output
projection.  Expand $A_x^a(-r-1)$ using
$\binom{-r-1}{m}=(-1)^m\binom{m+r}{m}$ to obtain
\eqref{cyquant:eq:quantum-kernel}.  Only $m\leq a$, $n\leq b$
contribute.  Every $A_x^a(m)$ kills the first vacuum, and
$R_y^a(r)\Psi_a(\mathbf1)=0$ for $r\geq0$; hence
\[
 W_u(w)(\mathbf1)=w,\qquad W_u(\mathbf1)(v)=e^{uT}v.
\]
Evaluation is the inverse onto the cyclic image, as in the classical
case.  Trace states remain, although their correction operators are zero.

A word of current indices $m_1,\ldots,m_d$ acting on weight at most
$b$ has intermediate weight at most
$b+1+\sum_i\max(0,-m_i)$ at this precision.  Choose $M$ above this
bound, and enlarge the ghost input block for its finite quadratic
rows.  Then the native action, its adjoint, and all cyclic comparisons
for the word stabilize together.  On input weights $i,j$ and output
weight $k$, the first-order Laurent exponent is $k-i-j-1$.
Thus the finite projected range is $[-a-b-1,c]$.
These bounds prove compatibility on common enlargements and with
every finite collection of derivatives.  Stabilization is asserted on
the cyclic image; no action on the entire Hom space is needed.
\end{proof}

The cancellation is visible on $v_E=J_E(-1)\mathbf1$ at the second
vacuum.  Each of the two sums in
\eqref{cyquant:eq:finite-right-dual} has $u^{-1}$ coefficient
$2Mv_E$ before subtraction: the zero-mode trace Casimir is $4$
on the adjoint representation.  Their difference is zero, as required
by $W_u(\mathbf1)(v_E)=e^{uT}v_E$.  Dropping the right-dual term
would leave a coefficient growing with the cutoff.
For $F=E_{21}$ and $H=E_{11}-E_{22}$, the first-order root coefficients
are
\[
 [u^{-3}\mathbf1](W_u-\Psi)(v_F)(v_E)=-\hbar K(K+2),
 \qquad
 [u^{-2}v_H](W_u-\Psi)(v_F)(v_E)=-\hbar(K+1).
\]
Both the scalar curvature contribution and the right-dual contribution
are needed; no denominator excludes zero or critical level.

The translation identities are
\[
 TW_u(w)-W_u(w)T=W_u(Tw),\qquad
 \partial_uW_u(w)(v)=W_u(w)(Tv).
\]
They identify the singleton coefficient diagrams as follows.
For finite PBW subsets $U\subset W$ with
$T\langle U\rangle\subset\langle W\rangle$, let $\iota_{UW}$ be
inclusion and $T_{WU}$ the finite translation matrix.  The perfect
coefficient transpose of the right-$\mathcal D$ Spencer complex gives
the pro diagram
\begin{equation}\label{cyquant:eq:ghost-spencer}
 \left\{
 \left[\mathcal D_{R_N}^{W^*}
       \xrightarrow{\ \partial\iota_{UW}^{\mathrm t}-T_{WU}^{\mathrm t}\ }
       \mathcal D_{R_N}^{U^*}\right]
 \right\}_{U,W},\qquad N=1,2,
\end{equation}
in degrees zero and one.  The right module convention is
$(f\,dz\otimes v)\partial=-f'\,dz\otimes v-f\,dz\otimes Tv$.
Formal transposition sends $\partial$ to $-\partial$, and the
singleton suspension supplies the displayed degrees.  Each map is
made before passing to the pro diagram; an ordinary inverse limit
of these complexes is not substituted for it.  The outgoing vacuum
map to $\omega[-1]$ is preserved.  It remains distinct from the
empty convolution unit.

The classical all-state comparison gave both ordered residue diagrams.
The first-order theorem above gives the native finite adjunction and
its singleton diagrams.  To identify the two quantum ordered products
one must still compare their separately typed tensor constraints,
both vacuum actions, arbitrary spectators, normal derivatives, and
the exchanged $13\mid2$ residue.  At precision three the native
adjunction must also incorporate \eqref{cyquant:eq:second-product},
\eqref{cyquant:eq:second-curvature}, and
\eqref{cyquant:eq:current-change}, reducing to this first-order map.
The all-mode singleton theorem does not replace these binary
requirements.

\section{The contact term of supported division}
\label{cyquant:sec:supported}

Let $a=z_2$, $w=z_1-z_2$, and $\kappa\ne0$.
Let $\widehat A_\kappa$ be the degreewise energy completion of the
commutative dg algebra on $x_n$ of cochain degree zero and $\xi_n$
of cochain degree one, for $n\geq0$, with
\[
 dx_n=\kappa(n+1)\xi_{n+1},\qquad d\xi_n=0,
 \qquad Tx_n=(n+1)x_{n+1},\quad T\xi_n=(n+1)\xi_{n+1}.
\]
Their energies are $n+1$ and $n$.  Extend scalars to $\mathbb C[a]$
before completing, so each energy coefficient remains polynomial in
$a$.  A first subscript $i$ denotes the copy at $z_i$.
On the disjoint locus, put $\theta=dw/w$ and place differential
forms before field coefficients.  In the total de Rham complex the
charge generators $y_i=\xi_{i,0}+\kappa^{-1}dz_i\,x_{i,0}$ are closed.
Let $\pi(y_i)=\epsilon_i$, where $\epsilon_i$ are the odd generators
of the exterior charge algebra.
The charge comparison with
$F_2(\epsilon_1,\epsilon_2)=0$ and
$F_2(\epsilon_2,\epsilon_1)=-\kappa^{-1}\theta$
differs from the binary roof representative
\[
 \mathcal C_2(y_1,y_2)=\frac{da}{\kappa w},\qquad
 \mathcal C_2(y_2,y_1)=-\frac{da+dw}{\kappa w}
\]
by opposite values $\beta,-\beta$, where $\beta=da/(\kappa w)$.
Their sum vanishes on commuting charge insertions, but
\[
 d\beta=-\frac{dw\wedge da}{\kappa w^2}.
\]
These two values do not define a homotopy on the whole bar construction.
Extending their division by $w$ to a supported coefficient requires
its actual polynomial module action.

Fix a finite-dimensional space $V$ in cochain degree zero and commuting endomorphisms
$A,B$.  Set $\Sigma=A+B$ and
$P_\Sigma=(\widehat A_\kappa\otimes V,d+\xi_0\Sigma)$.
For $e_n=[w^{-n}]$, $n\geq1$, and $e_0=0$, the supported carrier is
\begin{equation}\label{cyquant:eq:supported-carrier}
 \mathcal M_{A,B}=\bigcup_{r\geq1}\bigoplus_{n=1}^r e_nP_\Sigma,
 \qquad we_n=e_{n-1}.
\end{equation}
The field energies are completed at each fixed pole bound.  Each
vector has one finite pole bound independent of its energy; the
union and the energy product are not interchanged.

In its original normal frame, the differential and connections are
\[
 \begin{aligned}
 D(w)&=d+c(w)A+\xi_0B,&c(w)&=\sum_{n\geq0}w^n\xi_n,\\
 \nabla_w&=\partial_w+\kappa^{-1}j(w)A,&j(w)&=\sum_{n\geq0}w^nx_n,\\
 \nabla_a&=\partial_a+T+\kappa^{-1}(j(w)A+x_0B).
 \end{aligned}
\]
Only finitely many Taylor coefficients act at each pole bound.
The normal gauge
\begin{equation}\label{cyquant:eq:normal-gauge}
 \Gamma_A(w)=\exp\left(\frac A\kappa
                  \sum_{n\geq0}\frac{w^{n+1}x_n}{n+1}\right)
            =\sum_{n\geq0}w^n\gamma_n
\end{equation}
has $\gamma_0=1$ and energy $n$ in $\gamma_n$.
Multiplication by it intertwines these operators with
$d+\xi_0\Sigma$, $\partial_w$, and
$\partial_a+T+\kappa^{-1}x_0\Sigma$.

In the normal frame let $S(e_nm)=e_{n+1}m$ and define
\[
 S_A=\Gamma_A^{-1}S\Gamma_A,
 \qquad R_A(s)=[e_1](\Gamma_As),\qquad \iota(m)=e_1m.
\]
The bracket $[e_1]$ denotes coefficient extraction, equivalently the
normalized covariant residue of $dw\wedge s$.
The shift uses the unique negative-power representative of a principal
part.  It is not multiplication by a Laurent function on the quotient.

\begin{proposition}\label{cyquant:prop:contact}
The exact formulas and identities are
\begin{equation}\label{cyquant:eq:contact}
 \begin{gathered}
 S_A(e_nm)=e_{n+1}m-e_1\gamma_nm,\qquad
 R_A(e_nm)=\gamma_{n-1}m,\\
 wS_A=1,\qquad S_Aw=1-\iota R_A,
 \qquad [w,S_A]=\iota R_A,\qquad R_A\iota=1,\\
 [D(w),S_A]=0,\qquad [\nabla_a,S_A]=0,
 \qquad [\nabla_w,S_A]=-S_A^2.
 \end{gathered}
\end{equation}
The contact operator $\iota R_A$ is an idempotent cochain map.
The shift is continuous at each pole bound through all energies,
raises that bound by one, and does not lower energy.
\end{proposition}

\begin{proof}
On $e_nm$, the gauge retains precisely its first $n$ coefficients.
Therefore
$S\Gamma_A(e_nm)=\Gamma_A(e_{n+1}m)-e_1\gamma_nm$.
Since $\Gamma_A^{-1}(e_1v)=e_1v$, this proves both explicit formulas.
At pole bound $r$ there are only $r$ gauge coefficients; at a fixed
output energy each uses finitely many input energies.  This proves
continuity in the stated order of completion.

In the normal frame $wS=1$ and $Sw=1-e_1[e_1]$.
Conjugate by $\Gamma_A$, which commutes with $w$, to obtain the
contact identities and idempotence.  The normal differential and
tangential connection have no $w$-dependence, so they commute with
$S$.  Finally $\partial_we_n=-ne_{n+1}$ gives
$[\partial_w,S]e_n=-e_{n+2}=-S^2e_n$.  Conjugation proves all the
connection identities and cochain compatibility.
\end{proof}

For instance,
$\gamma_1=\kappa^{-1}Ax_0$ and
$\gamma_2=\frac12(\kappa^{-1}Ax_1+\kappa^{-2}A^2x_0^2)$.
Thus the differential-compatible shift already has a field correction
on $e_1v$; its commutator with $w$ there is still $e_1v$.

\begin{theorem}\label{cyquant:thm:contact-obstruction}
If $\ker(A+B)\ne0$, then $\iota R_A$ is not null-homotopic for
the internal differential, even after forgetting the module actions.
Consequently the cochain-homotopy class of $S_A$ contains no
$\mathbb C[a,w]$-linear cochain map and cannot be the linear term
of an $A_\infty$ endomorphism of that module.  The obstruction
persists on a quasi-isomorphic resolution with the corresponding
module action and maps.  It applies to $B=-A$ without a nilpotence
assumption on $A$.
\end{theorem}

\begin{proof}
For $0\ne v\in\ker(A+B)$, the vector $e_1v$ is closed: positive
powers of $w$ kill it and the remaining differential is
$e_1\xi_0(A+B)v=0$.  The complex has no negative cochain degrees,
so this vector is not a boundary.  The contact operator fixes it.
Hence its cohomology action is nonzero.

If $K-S_A=Dh+hD$, then $[D,w]=0$ gives
$[w,K]-[w,S_A]=D[w,h]+[w,h]D$.
A module-linear $K$ has $[w,K]=0$, contradicting the nonzero
cohomology action of $[w,S_A]$.  For an $A_\infty$ module map,
the identity with one algebra input $w$ similarly says that the
commutator of its linear term with $w$ is the differential of its
next component.  The same contradiction applies.  Compatible
quasi-isomorphisms conjugate the cohomology maps and cannot remove
this nonzero action.
\end{proof}

This is an obstruction to the specified class of supported division.
The two Laurent values $\beta,-\beta$ do not force an arbitrary
relative comparison homotopy into that class.  The remaining binary
problem is to solve
\[
 d_{\mathrm{rel}}H=\mathcal C-F\pi
\]
in the complex of continuous module cochains retaining finite Taylor
action and residue, or to prove that the actual right-hand side has a
nonzero class there.  Its other bar components and polynomial
representatives must be allowed before asserting a universal
obstruction.  A solution must then be compatible with every arity,
both ordered charts, spectators, all finite normal derivatives,
vacuum insertion, and the normal gauge.  The transverse arity and
the carrier's pole bound remain independent parameters throughout.
